% concatenated from article.tex, ex_shared.tex
% SIAM Article Template
\documentclass[hidelinks,onefignum,onetabnum]{siamart250211}

% Information that is shared between the article and the supplement
% (title and author information, macros, packages, etc.) goes into
% ex_shared.tex. If there is no supplement, this file can be included
% directly.

% SIAM Shared Information Template
% This is information that is shared between the main document and any
% supplement. If no supplement is required, then this information can
% be included directly in the main document.


% Packages and macros go here
\usepackage{lipsum}
\usepackage{amsfonts}
\usepackage{graphicx}
\usepackage{epstopdf}
\usepackage{algorithmic}
\ifpdf
  \DeclareGraphicsExtensions{.eps,.pdf,.png,.jpg}
\else
  \DeclareGraphicsExtensions{.eps}
\fi

% Add a serial/Oxford comma by default.
\newcommand{\creflastconjunction}{, and~}

% Used for creating new theorem and remark environments
\newsiamremark{remark}{Remark}
\newsiamremark{hypothesis}{Hypothesis}
\crefname{hypothesis}{Hypothesis}{Hypotheses}
\newsiamthm{claim}{Claim}
\newsiamremark{fact}{Fact}
\crefname{fact}{Fact}{Facts}

% Sets running headers as well as PDF title and authors
\headers{Conductance Estimation in Digraphs}{S. Shao, C. Yang, and X. Ye}

% Title. If the supplement option is on, then "Supplementary Material"
% is automatically inserted before the title.
\title{Conductance Estimation in Digraphs:  Submodular Transformation, Lov{\'a}sz Extension  and Dinkelbach Iteration\thanks{Submitted to the editors DATE.
\funding{This work was funded by the National Key R \& D Program of China (No.
2022YFA1005102) and the National Natural Science Foundation of China (Nos. 12325112, 12288101).}}}

% Authors: full names plus addresses.
\author{Sihong Shao\thanks{CAPT, LMAM and School of Mathematical Sciences, Peking University, Beijing 100871, China
  (\email{sihong@math.pku.edu.cn}).}% \url{http://www.imag.com/\string~ddoe/}).}
\and Chuan Yang\thanks{School of Mathematical Sciences, Peking University, Beijing 100871, China 
  (\email{chuanyang@pku.edu.cn}).}
\and Xinyang Ye\thanks{School of Mathematical Sciences, Peking University, Beijing 100871, China 
  (\email{xinyang@stu.pku.edu.cn}).}
}

\usepackage{amsopn}
\DeclareMathOperator{\diag}{diag}


%%% Local Variables: 
%%% mode:latex
%%% TeX-master: "ex_article"
%%% End: 


% Optional PDF information
\ifpdf
\hypersetup{
  pdftitle={Conductance Estimation in Digraphs:  Submodular Transformation, Lov{\'a}sz Extension  and Dinkelbach Iteration},
  pdfauthor={S. Shao and C. Yang and X. Ye}
}
\fi

% The next statement enables references to information in the
% supplement. See the xr-hyperref package for details.

%\externaldocument[][nocite]{ex_supplement}

% FundRef data to be entered by SIAM
%<funding-group specific-use="FundRef">
%<award-group>
%<funding-source>
%<named-content content-type="funder-name"> 
%</named-content> 
%<named-content content-type="funder-identifier"> 
%</named-content>
%</funding-source>
%<award-id> </award-id>
%</award-group>
%</funding-group>

\usepackage{subcaption}
\usepackage{booktabs}
\usepackage{amsmath}
\usepackage{amssymb}
\usepackage{cases}




\newtheorem{example}{Example}[section]
\renewcommand{\vec}[1]{\boldsymbol{\mathrm{#1}}}
%\renewcommand{\proofname}{\noindent \textit {Proof}}
\providecommand{\vw}{\ensuremath{\vec{w}}}
\providecommand{\vx}{\ensuremath{\vec{x}}}
\providecommand{\vs}{\ensuremath{\vec{s}}}
\providecommand{\vY}{\ensuremath{\vec{Y}}}
\providecommand{\vy}{\ensuremath{\vec{y}}}
\providecommand{\vd}{\ensuremath{\vec{d}}}
\providecommand{\vh}{\ensuremath{\vec{h}}}
\providecommand{\vp}{\ensuremath{\vec{p}}}
\providecommand{\vq}{\ensuremath{\vec{q}}}
\providecommand{\vl}{\ensuremath{\vec{l}}}
\providecommand{\va}{\ensuremath{\vec{a}}}
\providecommand{\vb}{\ensuremath{\vec{b}}}
\providecommand{\vu}{\ensuremath{\vec{u}}}
\providecommand{\vv}{\ensuremath{\vec{v}}}
\providecommand{\vy}{\ensuremath{\vec{y}}}
\providecommand{\vmu}{\ensuremath{\vec{\mu}}}
\providecommand{\cL}{\mathcal{L}}
\providecommand{\mX}{\mathbf{X}}
%\providecommand{\cut}{\mathbf{cut}}
\newcommand{\cM}{\mathcal{M}}
\newcommand{\cC}{\mathcal{C}}

\newcommand{\hw}{\hat{w}}
\newcommand{\R}{\mathbb{R}}
\newcommand{\1}{\vec{1}}
\newcommand{\ones}{\pmb{1}}
\renewcommand{\dot}[1]{\langle {#1} \rangle}
\newcommand{\norm}[1]{\lVert #1 \rVert}
\newcommand{\ind}{\mathrm{ind}}
\newcommand{\sign}{\mathsf{Sign}}
\newcommand{\sgn}{\mathsf{Sgn}}
\newcommand{\vol}{\mathrm{Vol}}
\newcommand{\supp}{\mathrm{supp}}
\newcommand{\cut}{\mathrm{cut}}


\newcommand{\bbR}{\mathbb{R}}
\newcommand{\bmd}{{\bm d}}
\newcommand{\bmx}{{\bm x}}
\newcommand{\bmw}{{\bm w}}
\newcommand{\bmz}{{\bm z}}
\newcommand{\bmone}{{\bm 1}}
\newcommand{\bmzero}{{\bm 0}}
\newcommand{\bmmu}{{\bm \mu}}
\newcommand{\bmchi}{{\bm \chi}}
\newcommand{\bmeta}{{\bm \eta}}
\newcommand{\caF}{\mathcal{F}}
\newcommand{\caL}{\mathcal{L}}
\newcommand{\caR}{\mathcal{R}}
\newcommand{\argmin}{\mathop{\mathrm{argmin}}}
\newcommand{\Lovasz}{Lov{\'a}sz\xspace}
\newcommand{\defeq}{\stackrel{\mathrm{\scriptscriptstyle def}}{=}}

\DeclareMathOperator*{\argmax}{\ensuremath{\mathrm{argmax\,}}}
\DeclareMathOperator*{\nen}{\ensuremath{\mathrm{NEN\,}}}
%\newtheorem*{proofsketch}{\noindent \textit {Proof}}

\begin{document}
\maketitle
\begin{abstract}
Conventional spectral digraph partitioning methods typically symmetrize the adjacency matrix, thereby transforming the directed graph partitioning problem into an undirected one, where bipartitioning is commonly linked to minimizing graph conductance. However, such symmetrization approaches disregard the directional dependencies of edges in digraphs, failing to capture the inherent imbalance crucial to directed network modeling. 
Building on the parallels between digraph conductance and conductance under submodular transformations, we develop a generalized framework to derive their continuous formulations. By leveraging properties of the Lov{\'a}sz extension, this framework addresses the fundamental asymmetry problem in digraph partitioning. We then formulate an equivalent fractional programming problem, relax it via a three-step Dinkelbach iteration procedure, and design the Directed Simple Iterative ($\mathbf{DSI}$) algorithm for estimating digraph conductance. The subproblem within $\mathbf{DSI}$ is analytically solvable, and the algorithm is guaranteed to converge provably to a binary local optimum. Extensive experiments on synthetic and real-world networks demonstrate that our $\mathbf{DSI}$ algorithm significantly outperforms several state-of-the-art methods in digraph conductance minimization.
\end{abstract}

% REQUIRED
\begin{keywords}
 digraph conductance, submodular transformation, fractional programming, Dinkelbach iterative framework, Lov{\'a}sz extension
\end{keywords}

% REQUIRED
\begin{MSCcodes}
05C20, 90C26, 90C27, 90C32, 90C30 % 68W25 % 15A42, 
\end{MSCcodes}






\section{Introduction}
\label{sec:intro}
Graph conductance is a crucial criterion to evaluate the quality of graph partitioning \cite{2006localPagerank,Benson2016HigherorderOO}, with applications in community detection and many other fields \cite{2000NomalizedCutImage,2013SpectralSparse}.
However, conductance is traditionally defined only for simple undirected networks, and algorithms for computing conductance in digraphs (directed graphs) have received less attention. Yet, many real-world networks are naturally represented as digraphs. In biological networks, digraphs signify carbon flow directions or connections between neurons \cite{Kaiser2006NonoptimalCP,Benson2016HigherorderOO}. In social media networks, they indicate the influence of podcast news on political parties or relationships between individuals \cite{2005politicalblog,2007emailcore,2017localhighorder}. In economic sciences, arcs represent transportation reachability between cities \cite{2007reachability}.



%\begin{remark}
%The term $ d^{\delta}_i$ cannot be entirely zero, as this would imply that the outdegree and indegree of each vertex are the same. This situation reduces to an undirected graph and is therefore not within the scope of consideration.
%\end{remark}
%To evaluate the quality of digraph partitioning, conductance $\varphi_D(S)$\eqref{eq:varphi_D} of a subset $S$ in a digraph was defined in \cite{2016nonlinearlap}.
For a weighted digraph $G=(V,A)$ and a partition $(S,\Bar{S})\, (\Bar{S} =V\backslash S, \,\,S \text{ and } \Bar{S}$ are both nonempty), where $V$ is the vertex set, $A$ is the arc set, the conductance $\varphi_D(S)$ of a subset $S$ in a digraph is  defined as% and $w_{ij}$ is the weight of each arc $i\to j$, 
\begin{equation}\label{eq:varphi_D}
\varphi_D(S)=\frac{\min\{\cut^+(S),\cut^-(S)\}}{\min\{\vol(S),\vol(\Bar{S})\}},
\end{equation} 
which is used to evaluate the quality of digraph partitioning \cite{2016nonlinearlap}.
% For a subset $S \subseteq V$, the out-boundary $\partial^+(S)$ and the in-boundary $\partial^-(S)$ are the set of arcs leaving and entering $S$, respectively. 
% $\partial^+(S) = \{(u, v) \in A \mid u \in S, v \in V \setminus S\}$ and 
% $\partial^-(S) = \{(u, v) \in A \mid u \in V \setminus S, v \in S\}$.
Here, $\cut^+(S)$ and $\cut^-(S)$ are the sum of weights of arcs leaving and entering $S$, respectively and $\cut(S) = \min\{\cut^+(S), \cut^-(S)\} .$ 
For an arc $e = i \to j,$ where $i$ is the arc starts from and $j$ is the arc ends to, denote the weight of this arc as $w_{ij}$ or $w_e.$
%For an arc $e = (u,v),$ denote it as $ u \to v,$ we say $u$ is the head and $v$ is the tail. $e$ is the arc leaving $u$ and is the arc entering $v.$
Set
 \begin{align*}
 \partial^+(S) = \{ i\to j \in A \mid i \in S, j \in V \setminus S\}, \text{ and }
 \partial^-(S) = \{ i\to j \in A \mid i \in V \setminus S, j \in S\}.
 \end{align*}
Then, $\cut^+(S) = \sum\limits_{e\in \partial^+(S)} w_e,$ $\cut^-(S) = \sum\limits_{e \in \partial^-(S)} w_e.$
The volume of $S$, denoted by $\vol(S)$, is defined as $\sum\limits_{v \in S} d_v,$ where $d_v$ is the sum of weights of all the arcs incident to vertex $v$ (containing both arcs leaving and entering $v$).
The out-degree $d^+_v$ and the in-degree $d^-_v$ of a vertex $v$ are defined as the sum of weights of arcs leaving and entering $v$, respectively. Then, the degree of $v$ is $d_v = d^+_v + d^-_v$ and the delta degree of $v$ is defined as $d_v^{\delta}=d^+_v-d^-_v $. To simplify notation, denote the corresponding degree vectors as $\vec d,$ $\vec d^+,$ $\vec d^-$ and $\vec d^{\delta}.$
In the same way, for a nonempty subset $ S\subsetneq V,$ there are out-conductance $\varphi^+(S)= \frac{\cut^+(S)}{\min\{\vol(S),\vol(\Bar{S})\}}$ and in-conductance $ \varphi^-(S)=\frac{\cut^-(S)}{\min\{\vol(S),\vol(\Bar{S})\}}$.
The conductance, in-conductance and out-conductance of a digraph $G$ are $\varphi_D(G) = \min\limits_{\emptyset \neq S \subsetneq V} \varphi_D(S),$ $\varphi^-(G) = \min\limits_{\emptyset \neq S \subsetneq V} \varphi^-(S),$ and $\varphi^+(G) = \min\limits_{\emptyset \neq S \subsetneq V}\varphi^+(S),$ respectively. 
In fact, it's easy to show that 
\begin{equation}\label{eq:min formula}
\varphi_D(G)=\min\{ \varphi^+(G),\varphi^-(G)\}.
\end{equation} We prove \eqref{eq:min formula} in Appendix \ref{proof of min formula}.

Many approaches have been proposed for spectral partitioning of digraphs. For example, the stationary distribution $\pi$ of the random walk process on digraphs can be used to define the Laplacian matrix \cite{Chung2005LaplaciansAT}, where the corresponding symmetrized adjacency matrix is $M = \frac{\Pi P + P^{T} \Pi}{2}$. Here, $P$ denotes the transition probability matrix of the random walk and $\Pi = \text{Diag}(\pi).$
Alternatively, there are several other symmetrization methods \cite{2011digraphSymmetic}, including $M = \frac{P + P^{T}}{2},$  $M=PP^T$ and $M = B + C,$ where 
$B = D_o^{-\alpha} P D_i^{-\beta} P^T D_o^{-\alpha}$, 
$C = D_i^{-\beta} P^T D_o^{-\alpha} P D_i^{-\beta},$ $D_o = \text{Diag}( \vec d^+),$ $D_i = \text{Diag}( \vec d^-),$ $\alpha,\beta$ are some parameters to be chosen.
However, these methods try to symmetrize the adjacency matrix of a digraph, since the eigenvector computation depends highly on the symmetry property of the adjacency matrix. 
These methods disregard the orientational dependencies of edges in digraphs, thus failing to model the intrinsic asymmetry of directed networks.
%These approaches ignore the directional dependencies of arcs in digraphs, thereby failing to capture the inherent imbalance of directed networks, a critical aspect of digraph modeling.
Hermite matrices were introduced in \cite{2020hermite_cluster} to represent adjacency relations in digraphs, thereby providing clustering algorithms for digraphs. However, this approach lacks theoretical guarantees.
A nonlinear Laplacian for digraphs and a heat kernel method for computing eigenvectors of the digraph adjacency matrix were developed in \cite{2016nonlinearlap}.
Through relaxation, it converts the conductance computing problem to the minimization of Rayleigh quotient $\mathcal{R}_G(\vx) = \frac{\vx^T \mathcal{L}_G(\vx) }{\vx^T \vx}$ for all $\vx\in \mathbb{R}^n \backslash \{0\}$ that are orthogonal to $\sqrt{\vec d/\vol(V)},$ where $\mathcal{L}_G$ is a nonlinear Laplacian operator.
However, this method needs to round the eigenvector to get the final partition, which is hard to obtain a local optimum.
%They also provided an algorithm to compute the Cheeger cut, demonstrating the convergence of the solution to the eigenvalue and eigenvector of the heat equation through finite iterations. 


Instead of using relaxation, we propose an equivalent formula of digraph conductance via the  Lov{\'a}sz extension.
We claim that
\begin{equation}\label{eq: main equivalent}
\min\limits_{\vx \in \mathbb{R}^n \backslash \{ t \vec 1\}, t \in R} r(\vx) =\min\limits_{\emptyset \neq S \subsetneq V } \varphi_D(S),
\end{equation}
%We define continuous function %$r(\vx)$ on $\R^n$ 
where
\begin{equation} 
\label{eq:dir ratio aim fun}
r(\vx)=\frac{\vol(V)\|\vx\|_\infty -I^+(\vx)-J(\vx)}{ 2N(\vx)},
\end{equation}
and
\begin{align}
I^+(\vx)&=\sum\limits_{i \to j \in A} {w_{ij}}|x_i+x_j|, \label{eq: Iplus}\\
J(\vx)&=|\sum\limits_{i \to j \in A}w_{ij}(x_i-x_j)|=|\sum\limits_{i \in V} d_{i}^\delta x_i |,\label{eq:Jx}\\
N(\vx)&= \min\limits_c \sum_{i\in V} d_i |x_i -c|.\label{eq:Nd}
\end{align}
%the optimum of function $r(\vx)$ is the same as the discrete aim function digraph conductance. 
Detailed proof of \eqref{eq: main equivalent} will be given in Section \ref{sec framework}. 
Motivated by the similarities between digraph conductance and conductance under submodular transformations, we propose a generalized framework to derive the equivalent continuous formulations.
%Given the similarities between digraph conductance and conductance under submodular transformations, we propose a generalized framework to derive equivalent continuous formulas for conductance under submodular transformations.
Details will be shown in Section \ref{sec framework} (Theorem \ref{them:min frac lovasz}).
%By \eqref{eq: main equivalent}, we could introduce continuous optimization methods to solve $\min\limits_{\vx \in \mathbb{R}^n \backslash \{ t \vec 1\}, t \in R} r(\vx),$ which has intensive study,such as the proximal gradient method \cite{2016ProximalGradientFP} and the Dinkelbach iterative framework \cite{1967Dinkelbach}. 

In this paper, we equivalently transform the fractional formula \eqref{eq: main equivalent} to a two-step iterative Dinkelbach  framework \eqref{eq:iterative alg two-step}, whose global optimal solution is equivalent to the minimum of \eqref{eq: main equivalent}.
However, the subproblem \eqref{eq:sub two-step} is NP-hard.
Through subgradient relaxation of \eqref{eq:sub two-step}, \eqref{eq:iterative alg two-step} can be modified to a three-step one \eqref{ite:three-step}.
%the two-step iterative framework \eqref{eq:iterative alg two-step} will converge to the optimum of \eqref{eq: main equivalent}.
%Our $\mathbf{DSI}$ (Directed Simple Iterative) method is aleveraging the properties of the Lov{\'a}sz extension \cite{Lovsz1982_SubmodularFA,2018Lovszsetpair} to transform the discrete optimization problem into a continuous one.
%For $\cut(S) = \min\{\cut^+(S), \cut^-({S})\}$, its Lov{\'a}sz extension $\cut^L(\vx)$ is not linear for each arc and is not additive. Direct calculation of its Lov{\'a}sz extension is not easy. So we choose an upper bound of the Lov{\'a}sz extension of $\cut(S)$ instead (see \eqref{eq: F(x)}).
Here, the subproblem \eqref{eq:subproblem} is analytically solvable and the sequence $\{r^k\}$ \eqref{eq:r sequence} is strictly decreasing with the special subgradient selection step (\eqref{eq；subg} and Algorithm \ref{alg:subgradient selection}). %will lead to the strictly decrease of $r(\vx).$ 
Based on the above properties, we show that our $\mathbf{DSI}$ algorithm will converge to a binary local optimum in finite steps.
%, since when the iteration step finishes, it will output a binary-valued vector through an analytical expression (Algorithm 2).
%Details of our $\mathbf{DSI}$ algorithm to solve $\min\limits_{\vx \in \mathbb{R}^n \backslash \{ t \vec 1\}, t \in R} r(\vx)$ will be given in Section 3.
%In fact, there is a large set 
%$$\mathcal{F}(\varphi_D) =\{ f:R^n \to \R| \min\limits_{\vx \in \mathbb{R}^n \backslash \{ t \vec 1\}, t \in R} f(\vx) =\min\limits_{\emptyset \neq S \subsetneq V } \varphi_D(S) \}.$$
%In fact, the selection of $r(\vx)$ is not unique. Let
%$$\mathcal{F}(\varphi_D) =\left\{ f:R^n \to \R| \min\limits_{\vx \in \mathbb{R}^n \backslash \{ t \vec 1\}, t \in R} f(\vx) =\min\limits_{\emptyset \neq S \subsetneq V } \varphi_D(S) \right\}.$$
%We choose suitable functions in $\mathcal{F}(\varphi_D)$ to improve the efficiency of solving process. 
We apply the $\mathbf{DSI}$ algorithm to various digraphs and demonstrate an improvement in conductance computation over some other methods.
In summary, our paper develops a simple and flexible method for computing digraph conductance with local convergence guarantee.
The $\mathbf{DSI}$ algorithm can also be generalized to estimate other conductance under submodular transformations (see Definition \ref{def:submodular transformations}).
By moving beyond the symmetric Laplacian of digraphs, our work opens new avenues for directed network analysis.






This paper is structured as follows.
%In Section 2, we will give some definitions for directed graph.
In Section \ref{sec framework}, we will prove \eqref{eq: main equivalent} and illustrate it in a more generalized framework utilizing submodular transformations.
%show that the steps to find elements in $\mathcal{F}$\ref{eq: set F} and illustrate this framework contains nearly all graph models.
In Section \ref{sec DSI}, we will present the detailed steps of our $\mathbf{DSI}$ algorithm for computing conductance in digraphs.
We will demonstrate that our algorithm exhibits a favorable convergence property (Theorem \ref{them:local convergence}), specifically converging to a local optimum within a finite number of steps.
Some proofs will be given in the Appendix for clarity.
In Section \ref{sec exp}, we will compare our $\mathbf{DSI}$ algorithm with other algorithms, including Sweepcut, on both synthetic dataset and real networks.
%and give some analysis on the numerical results.


\section{A General Framework for Seeking Equivalent Continuous Formulas}\label{sec framework}

The submodular transformation was proposed to unify and generalize the Cheeger inequalities for undirected graphs, directed graphs, and hypergraphs \cite{2019cheeger_submodular}. 
Under this framework, we derive continuous formulations of \eqref{eq:phi_F}, which inherently encompass the digraph conductance defined in \eqref{eq:varphi_D}.
%Since the digraph conductance \eqref{eq:varphi_D} and the conductance under transformations \eqref{eq:phi_F} share similarities, we take the digraph conductance as an example to design the algorithm.
%This Section mainly calculates the equivalent continuous function of $\varphi_F(S)$\eqref{eq:phi_F}.
We first give the definition of submodular transformation in Section \ref{subsec: submodular}.
Then, we introduce an important tool, Lov{\'a}sz extension in Section \ref{subsec: lovasz}.
Finally, we show how to find equivalent objective functions on multiple graph models in Section \ref{subsec: application}.
%5Then, we illustrate that our framework includes multiple cases of graph cut.


%For a submodular transformation $F,$ the cut function $\min\{\cut_F(S),\cut_F(V\backslash S)\}$ is not additive over each $e,$ and direct calculation of its Lov{\'a}sz extension is not easy.
%Instead of computing the Lov{\'a}sz extension itself, we identify other elements in $\mathcal{F}(\varphi_F)$ to design an algorithm. 


\subsection{Conductance under Submodular Transformations}\label{subsec: submodular}

Let $V$ be a set of size $n$ and $P(V) = \{S : S \subseteq V\}.$ $ f: P(V) \rightarrow \mathbb{R},$ is called \textbf{submodular} if for every $ S, T \subseteq V $, we have $ f(S) + f(T) \geq f(S \cap T) + f(S \cup T) $.
Note that the \textbf{cut function} $ \cut_G : \{0, 1\}^V \to \mathbb{R} $ associated with an undirected graph (a digraph) $ G $ is submodular, where $\cut_G(S) $ for a vertex set $ S \subseteq V $ represents the sum of weights of edges (arcs) leaving $ S $ and entering $ V \setminus S $. 
The submodular transformation was proposed by Yoshida \cite{2019cheeger_submodular} to give a general framework for graph conductance.
\begin{definition}[Submodular transformation \cite{2019cheeger_submodular}]\label{def:submodular transformations}
Let $G = (V, E)$ be a graph, $V$ and $E$ are its vertex set and edge set, respectively. Let $P(V) = \{S : S \subseteq V\},$ and $ F: P(V) \rightarrow \mathbb{R}^E$ is called \textbf{submodular transformation} if for each $ e \in E $, the induced component function $F_e: P(V) \to \mathbb{R}$ is submodular, where for any $S\subset V, $ $ F_e(S)$ is the corresponding component of $F(S)$ related to $e$.
%$ \forall e\in E,S \subseteq V$ .
\end{definition}
% and $f(S)(e)\geq 0,$$
%It's easy to check that 
%Through a submodular transformation, graphs can be transformed to digraphs, hypergraphs and other graph models.

\begin{example}[Digraph]
Let $G=(V,A)$ be a digraph. A submodular transformation $F:\{0,1\}^V\rightarrow\mathbb{R}^A$ relating to cut value can be obtained as follows. For each arc $e\in A$, let the corresponding component $F_e:P(V)\rightarrow\mathbb{R}$ denote the cut function of the digraph with the single arc $e$, where $F_e(S) =1$ if $e\in \cut^+(S)$ and  $F_e(S) =0$  otherwise. 
\end{example}
%Define a submodular transformation $F:\{0,1\}^V\rightarrow\mathbb{R}^A$, for each arc $e\in A$, $F_e:P(V)\rightarrow\mathbb{R}$ is the cut function of the digraph with a single arc $e$. To be more precise, for a subset $S\subsetneq V,$ $F(S)= \cut^+(S),$ and $F_e(S) =1_{e\in \cut^+(S)},$ $\forall e \in A.$
\begin{remark}
Note that the cut function $\cut_G(S)=\min\{ \cut^+(S),\cut^-(S)\}$ is not submodular. 
\end{remark}


%\begin{example}[Hypergraph]
%Let $G=(V,E)$ be a hypergraph. A submodular transformation $F:\{0,1\}^V\rightarrow\mathbb{R}^A$ relating to cut value can be obtained as follows. For each hyperedge $e\in E$, let the corresponding component $F_e:P(V)\rightarrow\mathbb{R}$ denote the cut function of the hypergraph with the single hyperedge $e$, which indicates $F_e(S) =1$ if $e\in \cut(S)$ and  $F_e(S) =0$  otherwise.
%$F_e:P(V)\rightarrow\mathbb{R}$ is the cut function of the hypergraph with a single hyperedge $e$. To be more precise, for a subset $S\subsetneq V,$ $F(S)= \cut(S),$ and $F_e(S) =1_{e\in \cut(S)},$ $\forall e \in E.$
%\end{example}


For a partition $(S,\Bar{S})$ of $V$ and a submodular transformation $F,$ the conductance of $S$ is defined as:
\begin{equation}
\label{eq:phi_F}
 \varphi_F(S) = \frac{\min\{\cut_F(S), \cut_F(V \backslash S)\}}{\min\{\vol(S), \vol(V \backslash S)\}},
 \end{equation}
%and the conductance $\varphi_F$ is the minimum over all $S,$ that is 
where $\cut_F(S) = \sum\limits_{e\in E} F_e(S),$
and $\varphi_F(G) = \min\limits_{\emptyset \neq S \subsetneq V} \varphi_F(S).$





\subsection{Lov{\'a}sz Extension and Beyond}\label{subsec: lovasz}

The main technique we use to find a continuous equivalent objective function of $\varphi_F(G)$ 
%for any $\emptyset \neq S \subsetneq V$ 
is Lov{\'a}sz extension \cite{Lovsz1982_SubmodularFA,Bac2013_learn_submodular_fun}, which
 is an important tool in establishing a bridge between discrete optimization and continuous optimization, thereby allowing the application of continuous optimization methods to solve discrete optimization problems \cite{2013_constrained_fractional,2018Lovszsetpair,2024maxcutSI}.


\begin{definition}[Lov{\'a}sz extension]\label{def:lovasz extension}
%Let $V = \{1, \ldots, n\} \subset \mathbb{N}$ 
For $\vx \in \mathbb{R}^n,$ let $ \sigma: [n]\cup \{0\} \to [n] \cup \{0\} $ be a bijection such that $x_{\sigma(1)} \leq x_{\sigma(2)} \leq \cdots \leq x_{\sigma(n)} $ and $ \sigma(0) = 0,$ where $ x_0 := 0$ and $[n]=\{1,2,...,n\}.$ 
Define the sets $V_{\sigma(i)} := \{j \in V : x_j > x_{\sigma(i)}\},$ for $ i = 1, \ldots, n - 1,$ and $ V_0 = V.$
Let $V$ be a set of size $n$ and $P(V) = \{S : S \subseteq V\}.$ Given $ f: P(V) \rightarrow [0, +\infty),$ the Lov{\'a}sz extension $F^L$ of $f $ is a mapping from $ \mathbb{R}^n $ to $ \mathbb{R}$ defined by
$$F^L(\vx) = \sum_{i=0}^{n-1} (x_{\sigma(i+1)} - x_{\sigma(i)}) f(V_{\sigma(i)}).$$
\end{definition}
%In the rest of this paper, for a function $f: P(V) \rightarrow [0, +\infty),$ we will use $F^L$ to denote its Lov{\'a}sz extension. 
%The following proposition is helpful in the computation of Lov{\'a}sz extensions.




\begin{proposition}\label{prop:lovasz extension}[\cite{Bac2013_learn_submodular_fun}, Definition 3.1]
Let $V$ be a set of size $n$ and $P(V) = \{S : S \subseteq V\}.$ Given $ f: P(V) \rightarrow [0, +\infty),$ 
then the Lov{\'a}sz extension $F^L$ of $f$ can be calculated as the following:
\begin{equation}
F^L(\vx) = \int^{\max\limits_{1 \leq i \leq n} x_i } _{\min\limits_{1 \leq i \leq n} x_i } f(V_t(\vx)) \, dt + f(V)\min\limits_{1 \leq i \leq n} x_i,
\label{eq: Lov{a}sz extension prop}
\end{equation}
where $V_t(\vx)=\{i\in V: x_i>t\}$.
\end{proposition}



Let $\mathbf{1}_S\in \mathbb{R}^n$ be the indicator vector of the set $S$, that is, the $i$-th component of $\mathbf{1}_S$ is 1 if and only if $v_i \in S$. Hence, for any $S\subseteq V$, we have $f(S) = F^L(\mathbf{1}_S)$. 
We first show some basic properties of Lov{\'a}sz extension.


\begin{lemma}
\label{lem:min Lovasz}
Let $V$ be a set of size $n$ and $P(V) = \{S : S \subseteq V\}.$ Given $g,$ $h: P(V) \rightarrow [0, +\infty),$ and define $f(S) = \min\{g(S), h(S)\},$ $\forall \, S\subset V.$ $F^l(\vx),$ $G^L(\vx),$ $H^L(\vx)$ are Lov{\'a}sz extensions of $f,$ $g,$ $h,$ respectively. Then $F^L$ is upper bounded by the minimum of $G^L$ and $H^L,$ that is,
$$F^L(\vx) \leq \min\{G^L(\vx), H^L(\vx)\}.$$
\end{lemma}
\begin{proof}

%$\sigma$ and $V_{\Sigma}$ are 
From Definition \ref{def:lovasz extension}, 

\begin{align*}
F^L(\vx) &= \sum_{i=0}^{n-1} (x_{\sigma(i+1)} - x_{\sigma(i)}) f(V_{\sigma(i)})\\
&= \sum_{i=0}^{n-1} (x_{\sigma(i+1)} - x_{\sigma(i)}) \min\{g(V_{\sigma(i)}),h(V_{\sigma(i)})\}\\
&\leq \min\left\{ \sum_{i=0}^{n-1} (x_{\sigma(i+1)} - x_{\sigma(i)}) g(V_{\sigma(i)}),\sum_{i=0}^{n-1} (x_{\sigma(i+1)} - x_{\sigma(i)}) h(V_{\sigma(i)})\right\}\\
&=\min\{G^L(\vx), H^L(\vx)\}.
\end{align*}
\end{proof}




\begin{theorem}\label{them:min frac lovasz}
Let $V$ be a set of size $n$ and $P(V) = \{S : S \subseteq V\}.$ Given $f_1,$ $f_2,$ $g: P(V) \rightarrow [0, +\infty),$ $g(\emptyset)=g(V)=0,$ and $\forall \, \emptyset \neq S \subsetneq V,$ $g(S)>0.$ Define $f(S) = \min\{f_1(S), f_2(S)\},$ $\forall \, S\subset V.$
$F^L,$ $F^L_1,$ $F^L_2,$ $G^L$ are Lov{\'a}sz extensions of $f,$ $f_1,$ $f_2,$ $g,$ respectively.
%The Lov{\'a}sz extensions of $f_1, f_2,g $ are $F^L_1, F^L_2,G^L$, respectively.
%Let $f(S) = \min\{f_1(S), f_2(S)\}$, $\forall S\subset V,$ and $F^L$ is the Lov{\'a}sz extension of $f.$
Set $\Tilde{F}^L = \min\{F_1^L, F_2^L\},$ then for any set $ S\subset V $, $ F^L(\mathbf{1}_S) = \Tilde{F}^L(\mathbf{1}_S) $.
If $f_1(V)=f_2(V)=0,$ then it holds that
$$\min_{\emptyset \neq S \subsetneq V} \frac{f(S)}{g(S)} = \inf_{\vx \in \mathbb{R}^n \backslash \{ t \vec 1\}, t \in R} \frac{F^L(\vx)}{G^L(\vx)}
 = \inf_{\vx \in \mathbb{R}^n \backslash \{ t \vec 1\}, t \in R} \frac{\Tilde{F}^L(\vx)}{G^L(\vx)}.$$
\end{theorem}


\begin{proof}

For any set $S\subset V$, 
$$F^L(\mathbf{1}_S) = f(S)=\min\{f_1(S),f_2(S)\} =\min\{ F_1^L(\mathbf{1}_S),F_2^L(\mathbf{1}_S)\}=\Tilde{F}^L(\mathbf{1}_S),$$
which proves $F^L(\mathbf{1}_S) =\Tilde{F}^L(\mathbf{1}_S).$
For any $\vx \in \mathbb{R}^n \backslash \{t \vec 1\}, t \in R,$ the corresponding set $V_{\sigma(i)}$ $(1\leq i \leq n-1)$ is not empty and is a proper subset of $V,$ which leads to $g(V_{\sigma(i)})>0,$ $\forall \, 1\leq i \leq n-1$ and $G^L(\vx)>0.$
Given that $f(V)=\min\{f_1(V),f_2(V)\}= g(V)=0,$  for any $\vx \in \mathbb{R}^n \backslash \{t \vec 1\}, t \in R,$ we have
\begin{align*}
F^L(\vx)&=\sum_{i=1}^{n-1} f(V_{\sigma(i)})\left(x_{\sigma(i+1)}-x_{\sigma(i)}\right).\\
&=\sum_{i=1}^{n-1}\frac{f(V_{\sigma(i)})}{g(V_{\sigma(i)})}g(V_{\sigma(i)})\left(x_{\sigma(i+1)}-x_{\sigma(i)}\right) \\
&\geq\min_{j=1,...,n-1}\frac{f(V_{\sigma(j)})}{g(V_{\sigma(j)})}\left(\sum_{i=1}^{n-1}g(V_{\sigma(i)}) (x_{\sigma(i+1)}-x_{\sigma(i)})\right)\\
&=\min_{j=1,...,n-1}\frac{f(V_{\sigma(j)})}{g(V_{\sigma(j)})} G^L(\vx).
\end{align*}
By Lemma \ref{lem:min Lovasz},
\begin{equation*}
\frac{\Tilde{F}^L(\vx)}{G^L(\vx)}\geq 
\frac{F^L(\vx)}{G^L(\vx)}\geq \min\limits_{j=1,...,n-1} \frac{f(V_{\sigma(j)})}{g(V_{\sigma(j)})}.
\end{equation*}
On the one hand,
$$\inf_{\vx \in \mathbb{R}^n \backslash \{ t \vec 1\}, t \in R}\ \frac{\Tilde{F}^L(\vx)}{G^L(\vx)}\geq
\inf_{\vx \in \mathbb{R}^n \backslash \{ t \vec 1\}, t \in R}\frac{F^L(\vx)}{G^L(\vx)}
\geq\min\limits_{j=1,...,n-1} \frac{f(V_{\sigma(j)})}{g(V_{\sigma(j)})}
\geq\min_{\emptyset \neq S \subsetneq V}\frac{f(S)}{g(S)}. $$
On the other hand, we have $\Tilde{F}^L(\mathbf{1}_S)=F^L(\mathbf{1}_S)$ for any set $S\subset V$, thus leading to
$$\min_{\emptyset \neq S \subsetneq V}\frac{f(S)}{ g(S)}
=\min_{\emptyset \neq S \subsetneq V}\frac{F^L(\mathbf{1}_S)}{G^L(\mathbf{1}_S)}
=\min_{\emptyset \neq S \subsetneq V}\frac{\Tilde{F}^L(\mathbf{1}_S)}{G^L(\mathbf{1}_S)}
\geq\inf_{\vx \in \mathbb{R}^n \backslash \{ t \vec 1\}, t \in R}\frac{\Tilde{F}^L(\vx)}{G^L(\vx)}.$$
Therefore,
$$\min_{\emptyset \neq S \subsetneq V}\frac{f(S)}{ g(S)}
=\inf_{\vx \in \mathbb{R}^n \backslash \{ t \vec 1\}, t \in R}\frac{\Tilde{F}^L(\vx)}{G^L(\vx)}. $$
\end{proof}


\begin{remark}\label{remark: moreH}
In fact, for any function $ H : \mathbb{R}^n \rightarrow \mathbb{R}$ satisfying $ H(\mathbf{1}_S) = F^L(\mathbf{1}_S)$ for any $ S\subsetneq V$ and $H(\vx) \geq F^L(\vx)$ for any $\vx\in \mathbb{R}^n$, 
$$\inf_{\vx \in \mathbb{R}^n \backslash \{ t \vec 1\}, t \in R}\frac{H(\vx)}{G^L(\vx)}=\inf_{\vx \in \mathbb{R}^n \backslash \{ t \vec 1\}, t \in R}\frac{F(\vx)}{G^L(\vx)}=\min_{\emptyset \neq S \subsetneq V}\frac{f(S)}{ g(S)}$$
holds, where $f,$ $g,$ $F^L,$ $G^L$ are defined in Theorem \ref{them:min frac lovasz}.
\end{remark}



%The  Lov{\'a}sz  extension possesses numerous advantageous properties, one of which is that the optimal values of the discrete and continuous functions coincide.
%In fact, except for the Lov{\'a}sz  extension itself, there are other continuous objective functions whose optima coincide with the discrete one.
%To be more precise, let


Theorem \ref{them:min frac lovasz} provides steps to identify equivalent continuous objective functions of \eqref{eq:phi_F}, which are not unique by Remark \ref{remark: moreH}.  Define
\begin{equation}\label{eq: set F}
\mathcal{F}(\varphi_F)= \left\{ f:\mathbb{R}^n \to \R| \min\limits_{\vx \in \mathbb{R}^n \backslash \{ t \vec 1\}, t \in R} f(\vx) =\varphi_F(G) \right\}
\end{equation}
as a set of all the equivalent continuous optimizations for $\varphi_F(G)$.
%Elements in $\mathcal{F}(\varphi_F)$ is not unique and
%we provide steps to find elements in $\mathcal{F}(\varphi_F).$
%a large set of equivalent continuous objective functions of $\varphi_F(S)$, that is, we show that there is a nonempty set $\mathcal{F}(\varphi_F)$ such that 
%Moreover, different selection of $f\in\mathcal{F}(\varphi_F)$ will affect the solving efficiency, which will be illustrated in Section 3 in detail.








\begin{corollary} \label{eq；conclusion}
From Remark \ref{remark: moreH}, for any function $F(\vx)$ satisfying
\begin{equation}\label{eq:equal property}
F(\vx)\geq \min\{ \cut_F^{+L}(\vx),\cut_F^{-L}(\vx)\} \text{ and }  F(\mathbf{1}_S)=\min\{\cut^+_F(S),\cut^-_F(S)\}
\end{equation}
for any $\vx\in\mathbb{R}^n$ and $S\subset V$, the following equality holds:
 $$\varphi_F(G) = \min\limits_{\vx \in \mathbb{R}^n \backslash \{ t \vec 1\}, t \in R}\frac{F(\vx)}{N(\vx)},$$ 
%$\min\limits_{\emptyset \neq S \subsetneq V}\varphi_F(S)$ is equivalent to solve $\min\limits_{\vx \in \mathbb{R}^n \backslash \{ t \vec 1\}, t \in R}\frac{F(\vx)}{N(x)},$ 
where $N(\vx)$, $\cut^{+L}(\vx)$, $\cut^{-L}(\vx)$ are the Lov{\'a}sz extensions of $\min\{\vol(S), \vol(\Bar{S})\},$ $\cut^+(S)$, $\cut^-(S),$  respectively.
Therefore, we have $\frac{F(\vx)}{N(\vx)}\in\mathcal{F}(\varphi_F)$.
\end{corollary}

%\subsection{Submodulart Transformation Cases}
%In this subsection,

\subsection{Applications} \label{subsec: application}
We now provide more than one equivalent objective function of digraph conductance and undirected graph conductance using Corollary \ref{eq；conclusion}.

\paragraph{Directed Conductance}\label{para:digraph}
Now, we will use Theorem \ref{them:min frac lovasz} to explore equivalent continuous objective functions for $\varphi_D(S)$ \eqref{eq:varphi_D} and prove \eqref{eq: main equivalent}. 
%Let $G = (V, A)$ be a digraph, $V$ is the vertex set and $A$ is the arc set. The out-degree $d^+_v$ and the in-degree $d^-_v$ of a vertex $v$ are the sum of weights of arcs leaving and entering $v$, respectively. Then, the degree of $v$ is defined as $d_v = d^+_v + d^-_v$ and the delta degree is defined as $d_v^{\delta}=d^+_v-d^-_v $. 
%For a weighted digraph $G=(D,A),$ let $w_{ij}$ be the weight of arc $\{i,j\} \in A.$
First, we calculate the Lov{\'a}sz extension $\cut^{+L}(\vx)$ for $\cut^+(S)$.

\begin{lemma}
Let $G=(V,A)$ be a digraph, 
the Lov{\'a}sz extension of $\cut^+(S)$ is
$$\cut^{+L}(\vx) = \sum_{i \to j \in A} \frac{w_{ij}}{2}(x_i - x_j + |x_i - x_j|).$$
\end{lemma}

\begin{proof}
Given $\vx\in\mathbb{R}^n,$ and set $V_t(\vx) = \{i \in V : x_i > t\}$, the sum of the weights of the arcs leaving $V_t(\vx)$ is
$$\cut^+(V_t(\vx)) = \sum_{i \to j \in A} w_{ij} \chi_{x_j \leq t < x_i}.$$
Since $\cut^+(V) = 0,$ and by Proposition \ref{prop:lovasz extension}, we obtain
\begin{align*}
\cut^{+L}(\vx) &= \sum_{i \to j \in A} w_{ij} \int^{\max\limits_{1 \leq i \leq n} x_i } _{\min\limits_{1 \leq i \leq n} x_i } \chi_{x_j \leq t < x_i} dt \\
&= \sum_{i \to j \in A} w_{ij} \max\{0, x_i - x_j\} \\
&= \sum_{i \to j \in A} \frac{w_{ij}}{2}(x_i - x_j + |x_i - x_j|)
.
\end{align*}
%The last equality uses the identity transformation
%\begin{equation}\label{eq:max}
% \max\{a, b\} = \frac{1}{2} (a + b + |a - b|).
%\end{equation}
\end{proof}
Similarly, for $\cut^-({S})$, its Lov{\'a}sz extension is
$$\cut^{-L}(\vx) = \sum_{i \to j \in A} \frac{w_{ij}}{2}(x_j - x_i + |x_j - x_i|).$$
For $\cut(S) = \min\{\cut^+(S), \cut^-({S})\}$, its Lov{\'a}sz extension $\cut^L(\vx)$ is not linear for each arc and is not additive. Instead of directly calculating $\cut^L(\vx)$, we notice it can be upper bounded by $\cut^{+L}(\vx)$ and $\cut^{-L}(\vx).$ Under condition \eqref{eq:equal property}, the upper bound function will not violate the equivalent condition.
From Lemma \ref{lem:min Lovasz}, for all $\vx\in\mathbb{R}^n$, we have% we know that $\cut^L(\vx)$ can be upper bounded by $\min\{\cut^{+L}(\vx),\cut^{-L}(\vx)\},$ that is 
%$\min\{ \cut^{-L}, \cut^{+L}\},$ where $\cut^{+L}$ and $\cut^{-L}$ are the Lov{\'a}sz extensions of $\cut^+(S)$ and $\cut^-(S)$, respectively. 
%In order to make the solution can be solved analytically, we find upper bound of the function $ I_D(\vx) $ can be transformed as follows:

\begin{align}
\cut^L(\vx)&\leq\min\{\cut^{+L}(\vx),\cut^{-L}(\vx)\}=\frac{1}{2}(I(\vx)-J(\vx)):=G(\vx),\\
G(\vx)&\leq \frac{1}{2}\left(\vol(V)\|\vx\|_\infty -I^+(\vx) -J(\vx)\right):=F(\vx),\label{eq: F(x)}
\end{align}
where $I(\vx) =\sum\limits_{i \to j \in A} {w_{ij}}|x_i-x_j|.$ 
%The first equality used the expression that $ \min\{a, b\} = \frac{1}{2} (a + b - |a - b|),$ and 
%The second inequality used the fact that $|a-b|+|a+b|= 2\max\{|a|,|b|\}.$
%Note that $r(\vx)$ can be rewritten as $ \frac{F(\vx)}{N(\vx)}$ and set $\Tilde{r}(\vx)=\frac{G(\vx)}{N(\vx)}.$ 
It's easy to check that for any set $ S \subset V $, $ F(\mathbf{1}_S) = \min\{\cut^+(S),\cut^-(S) \}$,
%\begin{equation}\label{eq: r and rtilde}
%\widetilde{r}(\vx) = \frac{I_D(\vx)}{N(\vx)},\,\,
% r(\vx) = \frac{F(\vx)}{N(\vx)}.
%\end{equation}
thus according to Corollary \ref{eq；conclusion}, we have
\begin{equation*}
\min\limits_{\vx \in \mathbb{R}^n \backslash \{ t \vec 1\}, t \in R} \Tilde{r}(\vx) = \min\limits_{\vx \in \mathbb{R}^n \backslash \{ t \vec 1\}, t \in R} r(\vx) = \min\limits_{\emptyset \neq S \subsetneq V} \varphi_D(S),
\end{equation*}
where $r(\vx)=\frac{F(\vx)}{N(\vx)}$, $\Tilde{r}(\vx)=\frac{G(\vx)}{N(\vx)}$, and now
%$$\min\limits_{\vx \in \mathbb{R}^n \backslash \{ t \vec 1\}, t \in R}\frac{F(\vx)}{N(\vx)} = \min\limits_{\vx \in \mathbb{R}^n \backslash \{ t \vec 1\}, t \in R}\frac{I_D(\vx)}{N(\vx)} = \min\limits_{\emptyset \neq S \subsetneq V} \varphi_D(S),$$
finishes the proof of \eqref{eq: main equivalent}.




\paragraph{Single Directed Conductance}
Given a digraph $G=(V,A),$ the equivalent continuous formulations for the out-conductance $\varphi^+(G)$ mentioned in Section~\ref{sec:intro} are
$$\min\limits_{\vx \in \mathbb{R}^n \backslash \{ t \vec 1\}, t \in R}\frac{\vol(V)\|\vx\|_\infty -I^+(\vx)-J_0(\vx)}{ 2N(\vx)}=\min\limits_{\vx \in \mathbb{R}^n \backslash \{ t \vec 1\}, t \in R}\frac{I(\vx)-J_0(\vx)}{ 2N(\vx)}=\varphi^+(G),$$
where $ J_0(\vx)=\sum\limits_{i \to j \in A}w_{ij}(x_i-x_j)=\sum_{i} d_{i}^\delta x_i$. For $\varphi^-(G),$ it is nearly the same.



\paragraph{Undirected Conductance \cite{Shao2024SI_cheeger,Chang20151Lapalg}}
Given an undirected graph $G=(V,E)$, let $\{i,j\}\in E$ denote the edge connecting $i\in V$ and $j\in V$. For any nonempty set $S\subsetneq V$, 
let $\cut(S)$ denote the sum of the weights of all edges crossing the partition $(S,\bar{S})$. The conductance of $G=(V,E)$ is defined as
\begin{equation}
 \varphi(G)=\min\limits_{\emptyset \neq S \subsetneq V}\varphi(S),\quad \varphi(S) = \frac{\cut(S)}{\min\{\vol(S), \vol(V \backslash S)\}}. 
 \end{equation}
By Corollary \ref{eq；conclusion}, the two equivalent continuous formulations of $\varphi(G)$ are
\begin{equation}
\label{eq:undirected}
    \min\limits_{\vx \in \mathbb{R}^n \backslash \{ t \vec 1\}, t \in R}\frac{\vol(V)\|\vx\|_\infty -\sum\limits_{\{i,j\} \in E} {w_{ij}}|x_i+x_j|}{ N(\vx)}=\min\limits_{\vx \in \mathbb{R}^n \backslash \{ t \vec 1\}, t \in R}\frac{\sum\limits_{\{i,j\} \in E} {w_{ij}}|x_i-x_j|}{ N(\vx)}.
\end{equation}
Given the non-uniqueness of equivalent continuous objective functions in $\mathcal{F}(\varphi_F)$, the selection of $r \in \mathcal{F}(\varphi_F)$ critically impacts computational efficiency. As demonstrated in the numerical experiments of \cite{Shao2024SI_cheeger}, iterative algorithms based on the first equivalent formulation of \eqref{eq:undirected} achieve significantly faster convergence than those using the second one. Motivated by this efficiency hierarchy, we will design iterative algorithm by adopting
$\min\limits_{\vx \in \mathbb{R}^n \backslash \{ t \vec 1\}, t \in R} r(\vx)$ \eqref{eq: main equivalent} for solving digraph conductance $\varphi_D(G)$.


\section{A Simple Iterative Algorithm}\label{sec DSI}



In this section, we elaborate on the $\mathbf{DSI}$ algorithm, which employs a continuous optimization approach to solve the discrete optimization problem.
Notably, this strategy has been applied to various graph cut problems,  including Max Cut \cite{2024maxcutSI}, AntiCheeger Cut \cite{2021AntiCheeger}, and Cheeger Cut \cite{Shao2024SI_cheeger} in undirected graphs.
%In this Section, we illustrate the details of the $\mathbf{DSI}$ algorithm, which introduce continuous optimization method to solve discrete optimization problem.
%Introducing continuous optimization method to solve discrete optimization problem is also used on other graph cut problems, such as Max Cut \cite{2024maxcutSI}, AntiCheeger Cut \cite{2021AntiCheeger}, and Cheeger Cut \cite{Shao2024SI_cheeger} in undirected graphs.
%These objective functions are all submodular functions, so their Lov{\'a}sz extensions directly produce convex (difference of convex) continuous equivalent objective functions.
%In this paper, we provide a more general framework for finding equivalent continuous objective functions that can include multiple graph models.
%fails to provide a convex continuous objective function. 
From \eqref{eq: main equivalent}, we know that 
%\min\limits_{\vx \in \mathbb{R}^n \backslash \{ t \vec 1\}, t \in R} r(\vx) =
$$\min\limits_{\vx \in \mathbb{R}^n \backslash \{ t \vec 1\}, t \in R}{r}(\vx) = \min\limits_{\emptyset \neq S \subsetneq V} \varphi_D(S),$$
where $r(\vx)$ is defined in \eqref{eq:dir ratio aim fun}.
%where $\widetilde{r}(\vx)$ is defined in 
Dinkelbach iterative algorithm \cite{1967Dinkelbach} was designed to solve fractional programming and here we use it to solve $\min\limits_{\vx \in \mathbb{R}^n \backslash \{ t \vec 1\}, t \in R} r(\vx).$
For notational convenience, we define
$\Omega_p^1 =\{ \vx \in \mathbb{R}^n: \max\limits_i x_i +\min\limits_i x_i = 0,\, \|\vx\|_p=1 \}$ and $\Omega_p^2 =\{ \vx \in \mathbb{R}^n: \min \limits_i |x_i| = 0,\, \|\vx\|_p=1 \}.$
Set
\begin{equation}\label{eq:Q}
Q_{r}(\vx)= \frac{I^+(\vx)+ J(\vx)+2r N(\vx) }{ \vol(V) }.
\end{equation}
Note that $\Omega_p^1 \cup \Omega_p^2$ is a compact set, and 
from Dinkelbach's algorithm, %\cite{1967Dinkelbach},
the sequence $\{\vx^k\}$ generated by the iterative steps
\begin{subequations}
\label{eq:iterative alg two-step}
\begin{numcases}{}
\vx^{k+1}=\argmin\limits_{\vx\in\Omega_p^1 \cup \Omega_p^2} 
\{\|\vx\|_\infty -Q_{r^k}(\vx)\},~p\in[1,\boldsymbol{\infty}], \label{eq:sub two-step}\\
r^{k+1}= r(\vx^{k+1}).
\end{numcases}
\end{subequations}
will converge to the optimum of $r(\vx),$ whose value is denoted as $r_{\min}.$ %\eqref{eq:dir ratio aim fun},$

\begin{theorem}\label{them:convergence}
The sequence $\{r^{k}\}$ generated by \eqref{eq:iterative alg two-step} will converge to the global optimum $r_{\min}.$ %$=\min\limits_{\vx \in \mathbb{R}^n \backslash \{ t \vec 1\}, t \in R} r(\vx).$ 
\end{theorem}
%\begin{proof}
%The detailed proof is given in Appendix \ref{them:convergence proof}.
%\end{proof}
\begin{proof}

It indicates $r_{\min}=\varphi_D(G)$. The definition of $\vx^{k + 1}$ in \eqref{eq:sub two-step} implies
$$0=\|\vx^{k}\|_{\infty}- Q_{r^{k}}(\vx^{k})
\geq\|\vx^{k + 1}\|_{\infty}-Q_{r^{k}}(\vx^{k + 1}),$$
which means $$ r^k \geq \frac{\vol(V)\|\vx^{k+1}\|_\infty - I(\vx^{k+1})- J(\vx^{k+1})}{2 N(\vx^{k+1}) }= r^{k + 1}\geq r_{\min},\ \forall\, k\in\mathbb{N}^{+}.$$
Thus we have
$$\exists\, r^{*}\in[r_{\min},r^{1}]\text{ s.t. }\lim_{k\rightarrow+\infty}r^{k}=r^{*}.$$
It suffices to show that $r_{\min} \geq r^{*}$. 
We denote
$$f(r)=\min_{\vx \in \Omega_p^1 \cup \Omega_p^2}(\|\vx\|_{\infty}-Q_r(\vx)),$$
%$$f(r)=\min_{\vx\in R^n\setminus\{0\}}(\|\vx\|_{\infty}-Q_r(\vx)),$$
which must be continuous on $\mathbb{R}$ by the compactness of $\Omega_p^1\cup \Omega_p^2$. %$\Omega_1\cup\Omega_2$. 
Note that
\begin{align*}
f(r^k)&=\|\vx^{k + 1}\|_{\infty} - Q_{r^k}(\vx^{k + 1})\\
&= \| \vx^{k + 1}\|_{\infty} -\frac{I^+(\vx^{k+1})+ J(\vx^{k+1})+2r^k N(\vx^{k+1}) }{ \vol(V) }\\
&= \frac{2N(\vx^{k+1})}{\vol(V)} \left(\frac{\vol(V) \| \vx^{k + 1}\|_{\infty} -I^+(\vx^{k+1})- J(\vx^{k+1}) }{2N(\vx^{k+1})}-r^k \right)\\
&=\frac{2N(\vx^{k + 1})}{\vol(V)}\cdot(r^{k + 1}-r^k)\to 0\quad\text{as}\quad k\to+\infty,
\end{align*}
%$Q_{r^k}(\vx)= \frac{I^+(\vx)+ J(\vx)+2r^k N(\vx) }{ \vol(V) },$
we have
$$f(r^*)=\lim_{k\to+\infty}f(r^k)=0,$$
and thus
$$\|\vx\|_{\infty} - Q_{r^*}(\vx)\geq 0,\, \forall\,\vx \in \Omega_p^1 \cup \Omega_p^2. $$
Hence, $r_{\min}\geq r^*,$ which completes the proof.
\end{proof}





Note that $Q_r(\vx)$ is convex, so we can introduce the subgradient method to make a relaxation.
Let $ g : \mathbb{R}^n \to \mathbb{R}$ be a convex function, and $ \vx\in \mathbb{R}^n $ be a point in the domain $ \text{dom} (g),$ a vector $ \vs \in \mathbb{R}^n $ is called a \textbf{subgradient} of $ g $ at $ \vx $ if it satisfies 
\begin{equation}\label{def: subdiff}
g(\vy) \geq g(\vx) + \langle \vs,\vy - \vx \rangle, \,\, \forall \,\vy \in \text{dom}(g),
\end{equation}
where $\langle \vv,\vb \rangle $ is the inner product of vector $\vv$ and $\vb.$ %$= \sum_i $
%the following inequality for all $\vy \in \text{dom} f $: $$f(\vy) \geq f(\vx) + \vs^T (\vy - \vx).$$ %\quad \forall \vy \in \text{dom} f.$$
The subdifferential of $g$ at $\vx$ is a set consisting of all the vectors $\vy$ satisfying inequality \eqref{def: subdiff}, denoted as $\partial g$.
%It is established that every interior point of a convex function admits a subgradient.
Since $Q_{r}(\vx)$ is convex and homogeneous of degree one, then for all $\vy\in\mathbb{R}^n,$ 
\begin{equation}\label{eq:Q relaxation}
\langle \vy, \vs \rangle \leq Q_{r}(\vy),\,\,\forall\, \vs \in \partial Q_{r}(\vx).%\,\,\forall \vy.
\end{equation}
Plugging the relaxation \eqref{eq:Q relaxation} into the subproblem \eqref{eq:sub two-step}, we get a relaxed version of \eqref{eq:sub two-step}. %and we can give exact solution of $\vx^{k+1}$ in each step.
The domain $\Omega_p^1 \cup \Omega_p^2$ can be relaxed to $\{\vx:\|\vx\|_p=1\}.$
This modifies the two-step Dinkelbach iterative scheme into the following three-step one:
\begin{subequations}
\label{ite:three-step}
\begin{numcases}{}
\vx^{k+1}=\argmin\limits_{\| \vx\|_p =1} 
\{ \|\vx\|_\infty -\langle \vx,\vs^k \rangle\},~p\in[1,\boldsymbol{\infty}], \label{eq:subproblem}\\
r^{k+1}= r(\vx^{k+1}), \label{eq:r sequence}\\
\vs^{k+1}\in \partial Q_{r^{k+1}}(\vx^{k+1}). \label{eq；subg}
\end{numcases}
\end{subequations}
It has been pointed out that the subproblem \eqref{eq:subproblem}
 admits an analytical solution \cite{2024maxcutSI}, with $\vx^{k+1}$ remaining non-constant (see Algorithm \ref{alg: exact solution} and Remark \ref{remark: nonconstant}). 
%We can write down a solution to the inner subproblem \eqref{eq:subproblem} in an analytical manner and the $\vx^{k+1}$ is still nonconstant.
Moreover, we are able to prove that such an iterative scheme still keeps the monotonicity (see Theorem \ref{them:monotonicity}) and has local convergence (see Theorem \ref{them:local convergence}).






\begin{algorithm}[htb]
\caption{ $\mathbf{DSI}$ algorithm for digraph conductance}
\label{alg:DSI}
\begin{algorithmic}
\STATE{\bf{Input:} an initial vector $\vx^1$, iterative number $T$}
\STATE{Set $k=1, \vx^*=\vx^1,r^*=r^1=r(\vx^1)$}
\WHILE{$k \leq T$}
\STATE{$\vs^k \gets \partial Q_{r^k}(\vx^k)$ (Algorithm \ref{alg:subgradient selection})}
%\STATE{}$r_k \gets (\vol(V)\|\vx^k\|_\infty -I^+(\vx^k)-J(\vx^k))/2N(\vx^k)$
%\STATE{}$\vs_k \gets (\vec p^k+\vec l^k+2r^k\vec a^k)/\vol(V)$
\STATE{$\vx^{k+1}=\argmin\limits_{\|\vx\|_p=1} \{ \|\vx\|_\infty -\langle \vx,\vs^k \rangle\},~p\in[1,\boldsymbol{\infty}]$ (Algorithm \ref{alg: exact solution})}
%\STATE{}$r_k \gets (\vol(V)\|\vx^k\|_\infty -I^+(\vx^k)-J(\vx^k))/2N(\vx^k)$
\STATE{$r^{k+1} \gets (\vol(V)\|\vx^{k+1}\|_\infty -I^+(\vx^{k+1})-J(\vx^{k+1}))/2N(\vx^{k+1})$}
\IF{$r^{k+1}<r^*$}
%\STATE{}$\vx^{k+1} \gets \vx^k$
\STATE{$r^*\gets r^{k+1}$, $\vx^*\gets \vx^{k+1}$, $k \gets k+1$}
%\STATE{$\vx^*\gets \vx^{k+1}$}
%\STATE{$k \gets k+1$}
\ELSE 
\STATE{break}
\ENDIF
\ENDWHILE
\RETURN $\vx^*,r^*$
\end{algorithmic}
\end{algorithm}




\begin{theorem}\label{them:monotonicity}
The sequence $\{r^{k}\}$ generated by Algorithm \ref{alg:DSI} satisfies $r^{k+1} \leq r^{k}$. 
\end{theorem}
%\begin{proofsketch}
%The detailed proof is given in Appendix \ref{them:monotonicity proof}.
%\end{proofsketch}





\begin{proof}
Given the definitions of  $r^{k+1}$ and $Q_{r^{k}}(\vx^k)$, we deduce that
\begin{align*}
r^{k+1}-r^k &=\frac{\vol(V)\|\vx^{k+1}\|_{\infty}-I^+(\vx^{k+1})-J(\vx^{k+1})-2r^k N(\vx^{k+1}) }{2N(\vx^{k+1})} \\
&=\frac{\vol(V)\|\vx^{k+1}\|_{\infty}-\vol(V) Q_{r^k}(\vx^{k+1}) }{2N(\vx^{k+1})} \\
&=\frac{\vol(V)}{2N(\vx^{k+1})}(\|\vx^{k+1}\|_\infty-Q_{r^k}(\vx^{k+1})).
\end{align*}
By the definition of  $\vx^{k+1},$ it follows that
\begin{align*}
0=\|\vx^k\|_\infty- \langle \vx^k,\vs^k \rangle \geq \|\vx^{k+1}\|_\infty-\langle \vx^{k+1},\vs^{k} \rangle.
\end{align*}
Since $Q_{r^k}(\vx^k)$ is convex and homogeneous of degree one, we have 
\begin{equation*}
\|\vx^{k+1}\|_\infty \leq \langle \vx^{k+1}, \vs^k \rangle \leq Q_{r^k}(\vx^{k+1}).
\end{equation*}
Therefore, $r^{k+1}\leq r^k.$
\end{proof}





\subsection{Subgradient Selection}


Since $\partial Q_{r^k}(\vx^k)$ is an interval, and the selection of subgradient $\vs^k\in\partial Q_{r^k}(\vx^k)$ will affect the decreasing speed of $\{r^k\}$ in the iteration framework \eqref{ite:three-step}. This relation is illustrated in the following Lemma.

\begin{lemma}[Lemma 3.1 in \cite{2024maxcutSI}]
\label{lem:realation rsL and decrease}
Suppose $\vx^k$, $r^k$ and $\vs^k$ are generated by the iteration \eqref{ite:three-step}.
Then for $k\geq 1$, we always have $ \Vert \vs^k\Vert_1 \geq 1$. In particular, 

(1) $ \Vert \vs^k\Vert_1= 1$ if and only if $\vx^k/\|\vx^k\|_\infty 
\in \sgn(\vs^k),$ where \begin{equation}
\sgn(t)=\left\{\begin{array}{ll}
{ \{1\},} & {\text{ if } t> 0,} \\ 
{[-1,1],} & {\text{ if } t= 0,} \\ 
{\{-1\},} & {\text{ if } t<0}.\\
\end{array} \right.\\
\end{equation}
%and 

(2) $ \Vert \vs^k\Vert_1>1 $ if and only if $L(\vs^k)<0,$
where \begin{equation}\label{eq:L}
L(\vs) := \min_{\|\vx\|_p=1}\{ \|\vx\|_\infty -\langle \vx,\vs \rangle\},\,\, \vs\in\mathbb{R}^n. \end{equation}
In addition, if $\| \vs^k \|_1 > 1$, then $ r^{k+1}< r^k$. 
\end{lemma}




Now we illustrate how to choose an ideal subgradient to output a local optimum.
According to ~\eqref{eq:Q}, we have
\begin{align}\label{eq:subgradient expression}
\vs^k =\frac{1}{\vol(V)}(\vu^k+\vy^k+2r^k\vv^k) \in \partial Q_{r^k}(\vx^k),
\end{align}
%\begin{align}\label{eq:subgradient expression}
%\vs^k :=\frac{1}{\vol(V)}(\vu^k+\vy^k+2r^k\vv^k) \in \partial Q_{r^k}(\vx^k)=\frac{1}{\vol(V)} ( \partial I^+(\vx^k)+\partial J(\vx^k) +2r^k \partial N(\vx^k)) ,
%\end{align}
where $\vu^k\in\partial I^+(\vx^k),$ $ \vy^k \in \partial J(\vx^k),$ $ \vv^k\in\partial N(\vx^k).$
Lemma \ref{lem:realation rsL and decrease} shows that the ideal $\vs^k$ satisfies that $\|\vs^k\|_1 >1,$ which will lead to $ r^{k+1}<r^k.$ Therefore, we select the subgradient $\vs^k \in \partial Q_{r^k}(\vx^k)$ to make $\|\vs^k\|_1$ as large as possible. $\partial Q_{r^k}(\vx^k)$ is an interval, which derives from the fact that $\partial I^+(\vx^k)$, $\partial J(\vx^k)$ and $\partial N(\vx^k)$ all consist of intervals. Their calculations are detailed as follows.
%$\partial Q_{r^k}(\vx^k)$ is an interval, we first use vector $\vl^k,$ $\va^k$ and $\vp^k,$ $\vq^k$ to characterize the lower bound and upper bound of each component of $\partial J(\vx^k),$ $\partial N(\vx^k)$ and $\partial I^+(\vx^k),$ respectively.
%Subsequently, using the formula of $\vl^k,$ $\va^k$ and $\vp^k,$ $\vq^k$, we select $\vy^k\in J(\vx^k),$ $\vv^k \in N(\vx^k)$ and $\vu^k \in I^+(\vx^k)$ to form $\vs^k$ \eqref{eq:subgradient expression} %and continue the iteration step\eqref{ite:three-step}.
For convenience, the superscript $k$ is omitted for simplicity in the rest of this section if there is no ambiguity.	
%We first characterize $\vl,$ $\va$ and $\vp,$ $\vq$, and then consider the choice of $\vu,$ $\vy,$ $\vv,$ respectively (See Algorithm \ref{alg:subgradient selection}).\\



	
	
$\bullet$ Characterization of $\partial I^+(\vx)$
	
%A direct algebraic calculation gives 
\begin{equation}
(\partial I^+(\vx))_i = [p_i-q_i,p_i+q_i],
\end{equation}
where 
\begin{align}
p_i &= \sum_{j:i\to j \text{ or } j \to i} w_{ij}\sign(x_i+x_j)-q_i, \label{eq:p}\\
q_i &= \sum_{j\in \nen(i,\vx)} w_{ij}, \label{eq:q}\\
\sign(t)&=\left\{\begin{array}{ll}
{ 1,} & {\text{ if } t\geq 0,} \\ 
{-1,} & {\text{ if } t <0,}\\\end{array} \right.\\
\nen(i,\vx) &= \{j:i\to j \text{ or } j \to i \text{ and }x_i+x_j=0\}.
\end{align}
%\begin{equation}\sign(t)=\left\{\begin{array}{ll}
%{ 1,} & {\text{ if } t\geq 0,} \\ 
%{-1,} & {\text{ if } t <0.}\\\end{array} \right.\\
%\end{equation}

%Here $\nen(i,\vx)$ denotes the set of ``negative equal neighbors" of the $i$-th vertex $i$ on $\vx$.\\


$\bullet$ Characterization of $\partial J(\vx)$


\begin{equation*}
\partial J(\vx)_i=\left\{\begin{array}{ll}
{l_i,} & {\text { if } J\neq 0,} \\ 
{[-|d^\delta_i|,|d^\delta_i|],} & {\text { if } J=0,}\\
\end{array}\right.\quad l_i=d_i^\delta \sign(\sum_i d_i^\delta x_i).
\end{equation*}
%where $$l_i=d_i^\delta \sign(\sum_i d_i^\delta x_i).$$



$\bullet$ Characterization of $\partial N(\vx)$

Let
\begin{align}
\text{median}(\vx)&=\argmin_{c\in\mathbb{R}}\sum_{i\in V}d_i|x_i-c|,\label{eq:median}\\
S^{\pm}(\vx)&=\{i\in V \big| x_i=\pm \|\vx\|_{\infty}\}, \label{eq:x_category1}\\
S^<(\vx)&=\{i\in V\big | |x_i|< \|\vx\|_{\infty}\},\label{eq:x_category2}\\
S^\alpha(\vx)&=\{i\in V \big |x_i=\alpha\},\quad\alpha\in \text{median}(\vx).\label{eq:x_category_alpha}
\end{align}

Then for a fixed $\alpha,$ the $i$-th component of $(\partial N(\vx))_i$ satisfies 
\begin{equation}
(\partial N(\vx))_i = [a_{i}^L, a_{i}^R], 
\end{equation}
where 
\begin{align}
a_{i}^L&=\left\{\begin{array}{ll}
{\max\{A-B+d_i,-d_i\},} & {\text { if } i\in S^{\alpha}(\vx)\text{ and }|	S^\alpha(\vx)|\geq 2,} \\ 
{A,} & {\text { if } i\in S^{\alpha}(\vx)\text{ and }|	S^\alpha(\vx)|\leq 1,}\\
{d_i\sign(x_i-\alpha),} & {\text { if }i\in V\setminus S^{\alpha}(\vx),}
\end{array}
\right. \label{eq:a_bound_l}\\
a_{i}^R &= \left\{\begin{array}{ll}
{\min\{A+B-d_i,d_i\},} & {\text { if } i\in S^{\alpha}(\vx)\text{ and }|	S^\alpha(\vx)|\geq 2,} \\ 
{A,} & {\text { if } i\in S^{\alpha}(\vx)\text{ and }|	S^\alpha(\vx)|\leq 1,}\\
{d_i\sign(x_i-\alpha),} & {\text { if }i\in V\setminus S^{\alpha}(\vx),}
\end{array}
\right. \label{eq:a_bound_r}	 \\
A &=\sum_{x_i<\alpha}d_i-\sum_{x_i>\alpha}d_i,\quad
B = \sum_{x_i=\alpha}d_i. \label{eq:v_ab}
\end{align}




In order to search for a subgradient that ensures $\|\vs^k\|_1>1 $ on the boundary of $\partial Q_r(\vx)$, we introduce a boundary indicator $\vb=(b_1,\ldots,b_n)\in\mathbb{R}^n,$
\begin{align}\label{eq:b}
b_{i}&=\left\{\begin{array}{ll}
{p_{i} +l_i+2r a_i+\chi_i q_i,} & {\text { if } J\neq 0,} \\ 
{p_{i}+ {\chi}_i |d_i^\delta|+ 2 r a_i+\chi_i q_{i},} & \text { if } J=0,
\end{array}\right.
\end{align}
with
\begin{align}\label{hh}	
a_i&=\begin{cases}
d_i\sign(x_i-\alpha),&\text{if $i\in V\setminus S^{\alpha}(\vx)$},\\
a^L_i,&\text{if $i\in S^{\alpha}(\vx)\cap S^+(\vx)$ and $|S^\alpha(\vx)|\geq 2$},\\
a^R_i,&\text{if $i\in S^{\alpha}(\vx)\cap S^-(\vx)$ and $|S^\alpha(\vx)|\geq 2$},\\ 
\argmax\limits_{t\in\{a^L_i,a^R_i\}}\{|p_i+l_i+2rt|\}, &\text{if $i\in S^{\alpha}(\vx)\cap S^<(\vx)$ and $|S^\alpha(\vx)|\geq 2$ and $J\neq 0 $},\\
\argmax\limits_{t\in\{a^L_i,a^R_i\}}\{|p_i+2rt|\}, &\text{if $i\in S^{\alpha}(\vx)\cap S^<(\vx)$ and $|S^\alpha(\vx)|\geq 2$ and $J= 0 $},\\
A,&\text{if $i\in S^{\alpha}(\vx)$ and $|S^\alpha(\vx)|\leq 1$},
\end{cases}
\end{align}
and 
\begin{equation}\label{eq:chi}
\chi(i)=\left\{\begin{array}{ll}
{\mp 1,} & {\text { if } x_i\in S^{\pm}(\vx),} \\ 
%{\sign(p_i+l_i+2 r a_i)),} & {\text { if } x_i\in S^{<}(\vx).}\\
{\sign(p_i+l_i+2ra_i),} & {\text{ If } x_i\in S^{<}(\vx) \text{ and } J\neq0 ,}\\
{\sign (p_i+2ra_i),} & {\text{ If }x_i\in S^{<}(\vx) \text{ and } J=0 .}\\
\end{array} \right.\\
\end{equation}
Each $b_i$ sits on the boundary of $(\partial Q_r(\vx))_i$.
Let $\Sigma_{\vb}(\vx)$ denote a collection of permutations of $\{1,2,\ldots,n\}$ such that for any $\sigma\in\Sigma_{\vb}(\vx)$, it holds $|b_{\sigma(1)}|\leq |b_{\sigma(2)}|\leq \cdots\leq |b_{\sigma(n)}|$.
Consequently, for any $i\in V$, it can be easily verified that 
\begin{equation}
\label{eq:desire_v}
\exists\, 1 \leq i \leq n,\, \text{ s.t. } \frac{x_i}{\|\vx\|_{\infty}}\notin\sgn(s_i) \Rightarrow \frac{x_i}{\|\vx\|_{\infty}}\notin\sgn(b_i).
\end{equation}
The steps of subgradient selection are given in Algorithm \ref{alg:subgradient selection}.
Let 
\begin{equation}\label{eq:Vb}
V_{\vb}=\argmax\limits_{i}\{ b_i \chi_i:\,b_i \chi_i>0\}.
\end{equation}
The iteration will stop if $V_{\vb}=\emptyset$, and the iterative values \eqref{eq:r sequence} will strictly decrease when $V_{\vb}\neq\emptyset.$
%we assume that for any $i \in V_b,$ $b_i \chi_i >0,$ otherwise $L(\vs)$ \eqref{eq:L} is not negative and $r^k$ will not strictly decrease, which will stop the iteration.
%We show that this way of selection has good property.



\begin{theorem}[$\mathbf{DSI}$: strict descent guarantee]\label{them:subgradient select strict decrease}
$V_b\neq \emptyset$ if and only if  $ \exists\, \vs \in \partial Q_{r}(\vx) \text{~such that~} {L}(\vs) < 0.$ Moreover, if $V_b\neq \emptyset,$ then ${L}(\vs^k) < 0,$ leading to $r^{k+1}<r^k$ due to Lemma \ref{lem:realation rsL and decrease},
where $\vs^k$ is selected through Algorithm \ref{alg:subgradient selection} at the $k$-th iteration.
% then $\vx^k$ and $r^k$ are produced by $\mathbf{DSI}$ at the $k$-th iteration, 
\end{theorem}
\begin{proof}
On the one hand, suppose $V_b\neq \emptyset$ and select $i \in V_b$ randomly.
Then we can select a subgradient $\vs\in \partial Q_r(\vx)$ such that $s_i = b_i/\vol(V).$ (The details of the selection steps are given in Algorithm \ref{alg:subgradient selection}).
Consequently,
\begin{equation}
\left\{
\begin{aligned} 
&x_{i}s_{i}<0, &\text{if }i\in S^{\pm}(\vx),\\
&	|s_{i}|>0, &\text{if }i\in S^<(\vx),\\
\end{aligned} \right.
\end{equation}
which directly yields ${L}(\vs)<0$ through Lemma \ref{lem:realation rsL and decrease}.

%The definition of $b_i$ and $b_i \chi_i > 0$ deduce that $x_i/\|\vx\|_\infty \notin \sgn(s_i)$ and this leads to $L(\vs)<0$ by Lemma \ref{lem:realation rsL and decrease}.

On the other hand, suppose that we can select a subgradient $\vs\in \partial Q_{r}(\vx)$ satisfying ${L}(\vs)<0$. According to Lemma~ \ref{lem:realation rsL and decrease}, there exists an $i^*\in\{1,2,\ldots,n\}$ satisfying
\begin{equation}
\left\{
\begin{aligned}
&\frac{b_{i^*}}{\vol(V)}\leq s_{i^*}<0, &\text{if }i^*\in S^+(\vx^k),\\
&\frac{b_{i^*}}{\vol(V)}\geq s_{i^*}>0, &\text{if }i^*\in S^-(\vx^k),\\
&	|b_{i^*}|\geq |s_{i^*}|>0, &\text{if }i^*\in S^<(\vx^k),\\
\end{aligned}
\right.
\end{equation}
and thus $\max\limits_i b_i\chi_i >0,$ which proves $V_b \neq \emptyset.$ 

%Furthermore, if $V_b\neq \emptyset$ then we can select the desired $i^*\in V_b$ in Algorithm \ref{alg:subgradient selection}. 

\end{proof}


\begin{algorithm}[htb]
\caption{Subgradient selection}
\label{alg:subgradient selection}
\begin{algorithmic}
\STATE{\bf{Input:}  $G=(V,A),$ $\vp$  \eqref{eq:p}, $\vq $ \eqref{eq:q},  $\va$ \eqref{hh}, $\vec \chi$ \eqref{eq:chi}, $\sigma\in\Sigma_{\vb}(\vx)$, $V_b$ \eqref{eq:Vb}}
\STATE{ // select $\vu$}
\STATE{randomly select $i^*\in V_{\vb}$, $u_{i^*}=p_{i^*}+\sum\limits_{j\in \nen(i^*,\vx)}w_{i^*j}\chi(i^*)$}
%\STATE{$u_{i^*}=p_{i^*}+\sum\limits_{j\in \nen(i^*,\vx)}w_{i^*j}\chi(i^*)$}
\FOR {each neighbor $i\in \nen(i^*,\vx)$} 
\STATE{$u_i=p_i+w_{ii^*}\chi(i^*)+\sum\limits_{j\in \nen(i,\vx)\backslash\{i^*\}}w_{ij}\chi(\argmax\limits_{t\in\{i,j\}}\{\sigma^{-1}(t)\})$}
\ENDFOR
\FOR {each remaining vertex $i$}
\STATE{$u_i = p_i+\sum\limits_{j\in \nen(i,\vx)}w_{ij}\chi(\argmax_{t\in\{i,j\}}\{\sigma^{-1}(t)\})$}
\ENDFOR
\STATE{// select $\vv$}
\IF {$|S^\alpha(\vx)|\leq 1$}
\STATE{$v_i=\begin{cases}
a_i&\text{if }i\in V\backslash S^{\alpha}(\vx)\\
\sum\limits_{i:x_i<\alpha }d_i - \sum\limits_{i:x_i>\alpha} d_i&\text{if }i\in S^{\alpha}(\vx)
\end{cases}$}
\ELSIF {$|S^\alpha(\vx)|\geq 2$}
\STATE{$j^*=\left\{\begin{array}{ll}{i^*} & {\text { if } i^*\in S^{\alpha}(\vx)} \\ 
{\argmax\limits_{t\in S^{\alpha}(\vx)}\{\sigma^{-1}(t)\}} & {\text { otherwise}}
\end{array}\right.$
\STATE{}
$v_i=\begin{cases}
a_i&\text{if }i\in \{j^*\} \cup V\backslash S^{\alpha}(\vx)\\
\frac{A-a_{j^*}}{B-d_{j^*}}d_i&\text{if }i\in S^{\alpha}(\vx)\backslash \{j^*\}
\end{cases}$}
\ENDIF
\STATE{// select $\vy$}
\STATE{$\vy=\begin{cases}
  \chi_{i^*}\vd^\delta  &\text{ if }J=0\\
  \sign(J)\vd^\delta &\text { otherwise}
\end{cases}$}
\STATE{$\vs=(\vu + \vy +2r \vv)/\vol(V)$}
\RETURN $\vs$
\end{algorithmic}
\end{algorithm}






%\begin{proof}
%The necessity is evident and we only need to prove the sufficiency.
%\end{proof}





\subsection{Local Convergence}


The proposed boundary-detected subgradient selection steps for the $\mathbf{DSI}$ algorithm ensure that the sequence $\{r^k\}$ strictly decreases whenever possible as shown in Theorem~ \ref{them:subgradient select strict decrease}.
Furthermore, we are able to prove that the $\mathbf{DSI}$ algorithm converges to a local minimum $\vx^*$ in $C_r$ within finite iterations from a discrete perspective, where
\begin{align}
C_r &= \{\vx\in\mathbb{R}^n \big| r(\vec y)\geq r(\vx), \,\, \forall\,\vec y\in\{R_i\vx: i\in\{1,\ldots,n\}\}\}, \label{eq:CT}\\
(R_i\vx)_j& =
\begin{cases}
x_j, & j\ne i,\\
-x_j, & j=i.
\end{cases}\label{eq:Tialp}
\end{align}

\begin{theorem}[local convergence]
\label{them:local convergence} Suppose that the output of the $\mathbf{DSI}$ algorithm converges to $\vx^*$, then $\vx^*\in C_r$.
\end{theorem}
%\begin{proof}
%The detailed proof is given in Appendix \ref{them: local convergence proof}.
%\end{proof}



\begin{proof} %[Proof of Theorem \ref{them:local convergence}]
According to the solution selection procedure in Algorithm \ref{alg: exact solution}, 
$\mathbf{\mathbf{DSI}}$ must produce nonconstant binary valued solutions (Remark \ref{remark: nonconstant}). Namely, $\vec \vx^*/\|\vec \vx^*\|_\infty \in \{-1, 1\}^n$. 
For simplicity, the superscript $*$ is neglected in the proof. Suppose $\vx\notin C_r$, which means there exists $i\in V$ satisfying $r(R_i\vx)<r(\vx).$
%Since $\|R_i\vx\|_\infty=\|\vx\|_\infty,$ then
%\begin{align*}r(R_i \vx)-r(\vx)=\frac{\vol(V)}{2N(R_i \vx)}(\|\vx^{k+1}\|_\infty-Q_{r^k}(\vx^{k+1})) =\frac{\vol(V)}{2N(R_i \vx)}(Q_r(\vx)_\infty-Q_{r}(R_i\vx)).\end{align*}
Given $\|R_i\vx\|_\infty = \|\vx\|_\infty,$ we deduce that
\begin{equation}\label{eq:ass}
	r(R_i\vx)<r(\vx) \quad \Leftrightarrow \quad Q_{r}(\vx)-Q_{r}(R_i\vx)<0,\quad r=r(\vx).
\end{equation}	
Hence,
\begin{align}
I^+(\vx)-I^+(R_i\vx)=\pm\sum\limits_{j: j\to i\in A \text{ or } i\to j\in A} 2w_{ij}x_j,\,\,i\in S^\pm(\vx), \label{eq:delta_iplus} \\
\begin{aligned}
N(\vx)-N(R_i\vx)=-|\Delta-d_i x_i|+|\Delta+d_i x_i|=&\begin{cases}
2d_ix_i,&\text{if }\Delta> d_i\|\vx\|_{\infty},\\
2\Delta x_i,&\text{if }\Delta \le d_i\|\vx\|_{\infty},\\ 
-2d_i x_i,&\text{if }\Delta< -d_i\|\vx\|_{\infty},
\end{cases}
\end{aligned}
\end{align}
where $\Delta=\sum\limits_{j=1}^nd_jx_j-d_ix_i$. Thus, it can be derived that %\text{if }\Delta<- d_i||\vx||_{\infty}
\begin{equation}\label{eq:mmm}
\begin{cases}
1\in\text{median}(\vx)\text{ and }1\in\text{median}(R_i\vx),&\text{if }\Delta> d_i\|\vx\|_{\infty},\\
x_i\in\text{median}(\vx)\text{ and }-x_i\in\text{median}(R_i\vx),&\text{if }\Delta\leq d_i\|\vx\|_{\infty},\\
-1\in\text{median}(\vx)\text{ and }-1\in\text{median}(R_i\vx),&\text{if }\Delta<-d_i\|\vx\|_{\infty}.
\end{cases}
\end{equation}	
We also have
\begin{align*}
J(\vx)-J(R_i\vx)=-|\kappa- d_i^\delta x_i|+|\kappa+ d_i^\delta x_i|=&\begin{cases}
2d_i^\delta x_i,&\text{if }\kappa> d_i^\delta\|\vx\|_{\infty},\\
2\kappa x_i,&\text{if }\kappa \le d_i^\delta\|\vx\|_{\infty},\\ 
-2 d_i^\delta x_i,&\text{if }\kappa< -d_i^\delta\|\vx\|_{\infty},
\end{cases}
\end{align*}
where $\kappa=\sum\limits_{j=1}^nd_j^\delta x_j-d_i^\delta x_i$. 
Let $p_i\in(\partial I^+(\vx))_i$ with 
$$p_i=\sum\limits_{j: j\to i\in A \text{ or } i\to j\in A} 2w_{ij}z_{j},\,\,\,z_{j} =x_j/\|\vx\|_{\infty}\in \sgn(x_i+x_j).$$
Then we choose $a_i\in(\partial N(\vx))_i$ and $l_i \in (\partial J(\vx))_i$ with 
$$a_i =\begin{cases}
d_i,&\text{if }\Delta> d_i\|\vx\|_{\infty},\\
\Delta,&\text{if }\Delta \le d_i\|\vx\|_{\infty},\\ 
-d_i,&\text{if }\Delta< -d_i\|\vx\|_{\infty},
\end{cases} \text{ and }\, l_i =\begin{cases}
d_i^\delta, &\text{if }\kappa> d_i^\delta\|\vx\|_{\infty},\\
\kappa,&\text{if }\kappa \le d_i^\delta\|\vx\|_{\infty},\\ 
- d_i^\delta,&\text{if }\kappa< -d_i^\delta\|\vx\|_{\infty},
\end{cases}$$
respectively. Accordingly, combining ~\eqref{eq:subgradient expression} and \eqref{eq:ass} together % and \eqref{eq:mmm} together 
leads to 
\begin{equation}
\label{xisileq0}
x_is_i=\frac{x_i p_i+x_i l_i+ 2rx_ia_i}{\vol(V)}<0,
\end{equation}
which contradicts Lemma~ \ref{lem:realation rsL and decrease}.
\end{proof}


\section{Numerical Experiments}\label{sec exp}

We implement our $\mathbf{DSI}$ algorithm using the Julia language. To achieve better solution quality, we select the initial point $ \vx_1 $ in Algorithm \ref{alg:DSI} as the second smallest eigenvector of the nonlinear Laplacian proposed in \cite{2016nonlinearlap}. Unlike their approach of simply rounding the scaled eigenvector, our method is rounding-free and yields superior solution quality.
We compare our algorithm against other digraph conductance algorithms in minimizing $\varphi_D$ \eqref{eq:varphi_D}
%global ratio cut objectives 
across various application domains.
%All experiments were conducted on a standard PC.
Since many algorithms are designed to partition a digraph into $k$ parts for any arbitrary integer $k,$ such as the CLSZ algorithm proposed in \cite{2020hermite_cluster} and the iterative algorithm presented in \cite{2024itealg_digraph},
their experimental results turn out to be unsatisfactory when $k=2.$ 
The Sweepcut algorithm introduced in \cite{2016nonlinearlap} is based on nonlinear digraph Laplacian and is originally designed for partitioning a graph into two parts.
Therefore, we only present the experimental results of our $\mathbf{DSI}$ algorithm and the Sweepcut algorithm.% introduced in \cite{2016nonlinearlap}.



\subsection{Experimental Results On Synthetic Dataset}
This Section mainly presents the experimental results of the $\mathbf{DSI}$ algorithm on synthetic datasets compared with Sweepcut algorithm.
We choose the DSBM (Directed Stochastic Block Model \cite{2020hermite_cluster}) to generate datasets for experiments.
DSBM is an extension of the traditional SBM (Stochastic Block Model \cite{Holland1983SBM}), which is a random model used to generate graphs with real clustering structures. This model controls the characteristics of the generated digraph through the following parameters: $k,n,p,q$ and $\eta.$
Here $k$ is the number of clusters and we set $k=2,$ $n$ is the size of each block, and $p,q,\eta$ are parameters used to generate edges. 
%We study graphs generated from the directed stochastic block model (DSBM) defined by $k,n,p,q$ and $\eta.$
To be more precise, the set $ \mathcal{G}( n, p, q, \eta) $ consists of graphs generated as follows: every $ G \in \mathcal{G} $ is a digraph defined on vertex set $ V = \{1, \dots, N\} $, where $ N = 2 n $. These vertices belong to clusters $ C_1, C_{2} $, where $ |C_j| = n $ for $ j \in [2] $. For any pair of vertices $u$ and $ v$, if they belong to the same cluster, they are connected by an arc with probability $ p $; otherwise, they are connected with probability $ q $. Moreover, if $ u \in C_1$ and $ v \in C_2$ are connected, the direction of this arc is determined by $\eta$: the direction is set to be $ u \to v $ with probability $\eta$, and $ v \to u $ with probability $1-\eta$. The direction of an arc inside a cluster is chosen uniformly at random. 


%To highlight the advantage of the $\mathbf{DSI}$ algorithm in solving DSBM,
%Since spectral techniques perform better in the SBM for large $p$, our focus is to compare the performance of different algorithms when $p$ is close to the connectivity threshold $log(N)/N$ of a random $G(n,p)$ graph. 
We generate digraphs sampled from DSBM with $N=2000$ and different $p,$ $q,$ and $\eta.$ To be more precise, we choose $p=q \in [0.005,0.006,0.008,0.10],$ $\eta\in [0,0.05,0.1,0.15,0.2,0.25,0.3]$. 
The experimental results are shown in Figure \ref{fig:results on DSBM} with respect to $\eta.$
%and exact numerical values are listed in Appendix \ref{sec: detail results on DSBM}.
%\paragraph{Numerical Analysis}
%We perform experiments on graphs randomly generated from the DSBM with different values of $n$, $p = q,$ and $\eta$. 
%Our reported results are averaged over 5 independently generated graphs for every fixed parameter set, and 
Sweep is conductance computed by the Sweepcut algorithm \cite{2016nonlinearlap}.
$\mathbf{DSI}$ is the conductance computed by the $\mathbf{DSI}$ algorithm.
%and the former add subgradient selection strategy, so it produce solution with higher quality.
%Sweep is conductance computed by \cite{2016nonlinearlap}.
It is clear that for different $p,q$ and $\eta$, our $\mathbf{DSI}$ algorithm provides better solution quality.






\begin{figure}[htb]
 \centering
 \begin{subfigure}{0.45\textwidth}
 \includegraphics[width=\linewidth]{Fig/DSBM_0.005_v2.pdf}
 \caption{$p=q=0.005$}
 \end{subfigure}
 \hfill
 \begin{subfigure}{0.45\textwidth}
 \includegraphics[width=\linewidth]{Fig/DSBM_0.006_v2.pdf}
 \caption{$p=q=0.006$}
 \end{subfigure}

 \begin{subfigure}{0.45\textwidth}
 \includegraphics[width=\linewidth]{Fig/DSBM_0.008_v2.pdf}
 \caption{$p=q=0.008$}
 \end{subfigure}
 \hfill
 \begin{subfigure}{0.45\textwidth}
 \includegraphics[width=\linewidth]{Fig/DSBM_0.01_v2.pdf}
 \caption{$p=q=0.01$}
 \end{subfigure}
 \caption{ Experimental results on DSBM: Figure (a) (b) (c) (d) illustrate the conductance values output by the $\mathbf{DSI}$ algorithm and the Sweepcut algorithm when $p=q=0.005,$ $p=q=0.006,$ $p=q=0.008$ and $p=q=0.01,$ respectively.}
 \label{fig:results on DSBM}
\end{figure}







\paragraph{Compare with Gurobi}


In order to better test the performance of the $\mathbf{DSI}$ algorithm, we will compare the experimental results of the $\mathbf{DSI}$ algorithm with the large-scale optimization solver Gurobi. %First, we will explain how to solve the problem using the Gurobi solver, transforming the digraph conductance computation problem into a mixed integer programming problem.
Since Gurobi's binary vectors can only take values of 0 and 1, we model it as:
$$\varphi_D(G)=\min\limits_{\vx\in\{0,1\}^n\text{ and } \vx \neq \vec 0, \vx \neq \vec 1 } \frac{\sum\limits_{i \to j\in A} w_{ij}(x_i+x_j-2x_ix_j)}{\min\left\{\sum\limits_i^n d_i x_i,\sum\limits_i^n d_i (1-x_i)\right\}}.$$

%$$x \text{ is not constant vector and} x\in\{0,1\}^n,$$

\begin{table}[htb]
\centering
\caption{Numerical results of $\mathbf{DSI}$ and Gurobi on DSBM when $p=q=0.005$ and $n=2000.$}
\begin{tabular}{ccccc}
%\begin{tabular}{|c|c|c|c|c|}
\toprule
$\eta$ & \textbf{$\mathbf{DSI}$} & \textbf{time} &\textbf{Gurobi} & \textbf{time} \\
\midrule
%0 & 0.0537 & 11.548 & 0.04 & 180.0564 \\
0.05 &0.0223 & 16.238 & 0.0682 & 120.0342 \\
0.10 &0.0379 & 18.215 & 0.0455 & 120.053 \\
0.15 &0.0575 & 16.816 & 0.0957 & 120.0339 \\
0.20 &0.0657 & 16.447 & 0.1667 & 120.0206 \\
0.25 &0.073 & 16.713 & 0.2047 & 120.0224 \\
0.30 &0.0824 & 16.798 & ${\inf}$ & 120.1067 \\
\bottomrule
\end{tabular}

\label{tab:my_label}
\end{table}
We set the time limit for Gurobi as 120s, inf means that Gurobi fails to output a solution.
Numerical results indicate that the $\mathbf{DSI}$ algorithm not only provides better solution quality than Gurobi, but also less time consuming.

%Numerical results indicate that Gurobi provides better solution quality for small-scale graphs, but the solution time is longer. When the graph scale becomes larger, with more than 1000 vertices and an increasing number of arcs, the solution quality significantly differs from that of our $\mathbf{DSI}$ algorithm.






\subsection{Experimental Results on Real Networks}
 
This section mainly presents the experimental results %of the $\mathbf{DSI}$ algorithm 
on real datasets. We first introduce the background of these real networks.

\paragraph{Florida Bay \cite{Benson2016HigherorderOO,Leskovec2014SNAPD}}

The Florida Bay dataset contains an unweighted digraph representing the carbon exchange between $ 125$ different species in Florida Bay. The arcs correspond to the typical diets of species observed within the bay.

\paragraph{C.elegans Network \cite{Kaiser2006NonoptimalCP}}
The C.elegans brain neural network describes the connections between neurons in the C.elegans brain. It is represented by a digraph containing 130 vertices and 610 arcs, where nodes represent neurons and arcs represent synaptic connections between neurons. 



\paragraph{Telegram \cite{2020dataTelgram}} The telegram dataset is a pairwise influence network between $245$ Telegram channels with $8,912$ arcs.
 
\paragraph{Blog \cite{2005politicalblog}} The Blog dataset records $19,024$ arcs between $ 1,212$ political blogs from the 2004 US presidential election.  \\


Table \ref{tab:exp realnet} provides the parameters and experimental results of these real networks.
Here, $n$ is number of vertices, $m$ is the number of arcs, and $c_n, c_m$ are the number of vertices and arcs in the largest connected component.
$\mathbf{DSI}$ is the conductance computed by the $\mathbf{DSI}$ algorithm.
%for $p=1$， and the former add subgradient selection strategy, so it produce solution with higher quality.
Sweep is conductance computed by Sweepcut algorithm.
It is clear that on different dataset, our $\mathbf{DSI}$ algorithm provides better solution quality.
%To show the clustering result of our algorithm better, we give several heat map of adjacency matrix of the partition in Figure 3. It is clear that our algorithm tends to find denser structures in directed networks.






\begin{table}[h!]
\centering
\caption{Numerical results of $\mathbf{DSI}$ and Sweepcut on real networks.}
\begin{tabular}{lcccccccc}
\toprule
name & $n$ & $m$ & $c_n$ & $c_m$ & $\mathbf{DSI}$ & {time 1} & $\text{Sweep}$ & {time 2} \\
\midrule
Celegans & 131 & 764 & 109 & 637 & 0.0126 & 5.953 & 0.0301 & 0.739 \\
Florida & 128 & 2106 & 103 & 1579 & 0.0087 & 6.4 & 0.0133 & 0.727 \\
blog & 1222 & 19024 & 793 & 15783 & 0.0286 & 8.678 & 0.0299 & 2.901 \\
telegram & 245 & 8912 & 138 & 4478 & 0.0209 & 8.633 & 0.0245 & 0.877 \\
\bottomrule
\end{tabular}
\label{tab:exp realnet}
\end{table}







%\input{1 intro}
%\input{2 aimfun}
%\input{3 dirSI}
%\input{4 dataset}
\section{Conclusion and Outlook}

In this paper, we present an algorithm for estimating digraph conductance that addresses the fundamental asymmetry challenge inherent in directed networks. To overcome this challenge, we leverage properties of the Lov{\'a}sz extension to derive equivalent continuous formulations for both digraph conductance and conductance under submodular transformations, thereby bridging the discrete-continuous optimization gap. 
Building on this novel continuous formulation, we introduce the Directed Simple Iterative ($\mathbf{DSI}$) algorithm for digraphs with the Dinkelbach iterative framework. The outputs of $\mathbf{DSI}$ are rigorously proven to converge to binary local optima. 
Comprehensive evaluations on synthetic and real-world datasets demonstrate that $\mathbf{DSI}$ significantly outperforms benchmark methods including Sweepcut and Gurobi.
This work breaks through the limitations of symmetric Laplacian-based approaches for digraphs, opening new avenues for research in the analysis of directed complex networks. Notably, our framework demonstrates the potential for generalization to estimate hypergraph conductance and other related combinatorial metrics, highlighting its broad applicability.





\appendix \label{sec appendix}
\section{Exact solution of \eqref{eq:subproblem}}

The exact solutions for various choices of norm $p$ were presented in  \cite{2024maxcutSI}. Different choices of $p$ will not affect the numerical results too much. Therefore, in this paper, we only consider the case when norm $p=1$, and the specific steps are detailed in Algorithm \ref{alg: exact solution}.


\begin{algorithm}[htb]
\caption{Exact solution of subproblem \eqref{eq:subproblem}}
\label{alg: exact solution}
\begin{algorithmic}
\STATE{ \bf{Input:} $\vs^k$}
\STATE{Sort $|\vs^k|$ s.t. $|s_{\pi(1)}^k|\geq |s_{\pi(2)}^k| \geq\dots\geq |s_{\pi(n)}^k| \geq |s_{n+1}^k|=0$}
\STATE{ $A_m^k=\sum\limits_{j=1}^m(|s_{\pi(j)}^k|-|s_{\pi(m+1)}^k|)$, $m_0=\min\{m: A_m^k>1\}$, $m_1=\max\{m: A_{m-1}^k<1\}$ }
%\STATE{ $m_0=\min\{m: A_m^k>1\}$ }
%\STATE{ $m_1=\max\{m: A_{m-1}^k<1\}$ }
\IF{$\| \vec \vs^k \|_1 >1 $} 
\IF{$1\leq i \leq m_1$}
\STATE{ $z_{\pi(i)}=1$}
\ELSIF{$m_0 < i \leq n$}
\STATE{ $z_{\pi(i)}=0$}
\ELSIF{$m_1< i \leq m_0$}
\STATE{$z_{\pi(i)}\in [0,1]$}
\ENDIF
\STATE{ $x_i^{k+1}= \sign(s_i^k) z_i / \|\vec z\|_1$ for $i=1, 2, \dots, n$ }
\ELSE 
\STATE{$\vx^{k+1}\in \sgn(\vs^k)/n$}
%$\|\vx^{k+1}\|_1=1$, and $\forall\,i\in V, x_i^{k+1}s_i^{k}\geq 0$.
\ENDIF
\RETURN $\vx^{k+1}$
\end{algorithmic}
\end{algorithm}


\begin{remark}\label{remark: nonconstant}
%From Lemma \ref{lem:realation rsL and decrease}, we know that if $r^{k+1}<r^k,$ then $L(\vs^k)<0.$   
We illustrate that if $L(\vs^k)<0,$ then the next point output by Algorithm \ref{alg: exact solution}
 is  also a nonconstant vector.
 Let $L(\vx, \vs^k):= \|\vx\|_\infty - \langle \vx,\vs^k\rangle .$
 Note that every constant vector $\vx$  will lead to $L(\vx,\vs^k) \geq 0.$  %  $\|\vx\|_\infty -\langle \vx, \vs^k\rangle \geq  0.$ 
To be more precise, if $\vx = t\vec 1,$ for some $ t \in R,$ then  $I(\vx)=\|\vx\|_\infty  \vol(V),$ $J(\vx)=0,$  $N(\vx)=0$  and $Q_{r^k}(\vx)= \| \vx\|_\infty.$  Therefore,
\begin{equation}
L(\vx,\vs^k)=\|\vx\|_\infty - \langle \vx,\vs^k\rangle \geq \|\vx\|_\infty - Q_{r^k}(\vx)=0.
\end{equation}
Thus, if $L(\vs^k)<0,$  then $\vx^{k+1}$ of the iteration \eqref{ite:three-step} found by Algorithm \ref{alg: exact solution} is nonconstant and  this will lead to  $r^{k+1}<r^k.$
To be more precise,
\begin{equation}
0 >L(\vs^k) =L(\vx^{k+1},\vs^k) %= \|\vx^{k+1}\|_\infty - \langle \vx^{k+1},\vs^k\rangle
\geq \|\vx^{k+1}\|_\infty - Q_{r^k}(\vx^{k+1}) 
%= \|\vx\|_\infty -\frac{I(\vx^{k+1} +J(\vx^{k+1}) +2 r^k N(\vx^{k+1}))}{\vol(V)}
%=\frac{\vol(V)\|\vx\|_\infty - I(\vx^{k+1} -J(\vx^{k+1}) -2 r^k N(\vx^{k+1}))}{\vol(V)}
=\frac{2 r N(\vx^{k+1})}{\vol(V)} (r^{k+1} - r^k).
\end{equation}
Subsequently, we have $r^{k+1}< r^k.$
\end{remark}



%Therefore, $m_1 <n.$ (If $m_1=n,$ then $z=\vec 1.$ This will lead to   $\vx^{k+1}\in \sgn(\vx^k)/n$ and $L(s^k)=0,$ a contradiction.)
%Note that $m_1\geq 1,$ 



%\begin{lemma}
%The minimizer $\vec \vx^{k+1}$ satisfies \begin{equation}\label{eq:m}\frac{\langle |\vec \vx^{k+1}|,|\vec \vs^k|\rangle}{\|\vec \vx^{k+1}\|_{\infty}} \ge \sum_{i=1}^{m} |s_i^k|,\end{equation} where $m =m_0$ for $\| \vec \vs^k \|_1>1$, and $m=n$ for $\| \vec \vs^k \|_1=1$. \end{lemma} 
%\begin{proof}
%The solution construction can be readily achieved using Theorem 3.1 and Lemma 3.2 in \cite{2024maxcutSI} and Aq.~\eqref{eq:m} comes from Corollary 3.1 in \cite{2024maxcutSI}. Here we skip the details for saving space.
%\end{proof}


\section{Proof of \eqref{eq:min formula}}
\label{proof of min formula}
\begin{proof}
On the one hand,  for any $S \subsetneq V,$ $\varphi_D(S) \leq \min\{\varphi^+(S),\varphi^-(S)\}.$
Therefore, $ \varphi_D(G) =\min\limits_{\emptyset \neq S \subsetneq V} \varphi_D(S) \leq \min\limits_{\emptyset \neq S \subsetneq V} \{ \varphi^+(S),\varphi^-(S)\} = \min\{\varphi^+(G),\varphi^-(G)\}.$ 

On the other hand, assume that $\varphi_D(S_0)=\varphi_D(G)$.
Without loss of generality,
assume that $\cut^+(S_0)\geq \cut^-(S_0).$
Subsequently, $\varphi_D(G)=\varphi_D(S_0) = \varphi^-(S_0) \geq \min\limits_{\emptyset \neq S \subsetneq V} \varphi^-(S) \geq \min\limits_{\emptyset \neq S \subsetneq V}\{\varphi^-(S),\varphi^+(S)\}.$ Therefore, $ \varphi_D(G)=\min\{ \varphi^+(G),\varphi^-(G)\}.$
\end{proof}





\bibliographystyle{siamplain}
\bibliography{references}
\end{document}
